%% ============================================================================
%% Chapter 11: Asynchronous & Event-Driven APIs: Webhooks, SSE & WebSockets
%% Mastering APIs: From Zero to Production Guru
%% ============================================================================
\chapter{Asynchronous \& Event-Driven APIs: Webhooks, SSE \& WebSockets}
\label{ch:async_apis}

%% ----------------------------------------------------------------------------
%% Executive Summary
%% ----------------------------------------------------------------------------
Every chapter up to this point has treated HTTP as a telephone call: the
client dials, the server answers, the connection hangs up. That
request/response shape is the right default for the overwhelming majority of
API work --- but it structurally cannot express a large and growing class of
interactions in which \emph{the server knows something first}. A robot's
battery crosses a critical threshold; a payment settles three minutes after
checkout; a colleague edits the document you are viewing. In each case the
interesting event originates server-side, and the client wants to know about
it in milliseconds, not at the next polling interval. This chapter is about
the three industrial-strength answers to that problem --- \textbf{webhooks},
\textbf{Server-Sent Events (SSE)}, and \textbf{WebSockets} --- and about the
engineering discipline each one demands before it can survive production.

We begin with the interaction-pattern taxonomy, because the most expensive
mistake in this space is choosing the wrong pattern, not implementing the
right one badly. Request/response, short polling, long polling, and server
push are compared on the axes that actually decide architectures: latency,
fan-out cost, infrastructure compatibility, cacheability, and failure
semantics. Webhooks are then treated with the seriousness they deserve as
\emph{reverse APIs}: you, the provider, become an HTTP client calling
endpoints you do not operate, across networks you do not control, with
delivery guarantees you must engineer explicitly. We specify the de-facto
industry contract wire by wire --- at-least-once delivery, bounded retries
with exponential backoff, HMAC-SHA256 signatures over timestamp-plus-body in
the Stripe-style \texttt{t=...,v1=...} header, timestamp tolerance windows
for replay protection, endpoint verification handshakes, secret rotation,
idempotent receivers, and dead-letter queues --- and we give the provider and
consumer checklists that turn that contract into an operational habit.

Server-Sent Events get the same wire-level treatment: the exact
\texttt{text/event-stream} grammar (the \texttt{event}, \texttt{data},
\texttt{id}, and \texttt{retry} fields, comment keep-alives, multi-line
\texttt{data} folding), the \texttt{EventSource} reconnection protocol with
\texttt{Last-Event-ID}, and the two production killers that rarely appear in
tutorials --- buffering intermediaries that swallow your stream
(\texttt{X-Accel-Buffering: no}) and the HTTP/1.1 six-connections-per-origin
limit that HTTP/2 removes. WebSockets are specified per RFC~6455 with the
full opening handshake --- including a worked derivation of
\texttt{Sec-WebSocket-Accept} as
$\mathrm{Base64}\bigl(\mathrm{SHA\text{-}1}(key \,\|\, \mathrm{GUID})\bigr)$ ---
the frame format byte by byte, masking and the cache-poisoning rationale for
it, opcodes, ping/pong keep-alive, close codes, and subprotocol negotiation.
Just as importantly, we enumerate when WebSockets is the \emph{wrong} tool:
stateless CRUD, cacheable reads, and anything that must traverse hostile
intermediaries. The conceptual frame throughout is the event-driven
architecture canon~\cite{kleppmann2017ddia} and its microservice
application~\cite{newman2021microservices}, with HTTP semantics as specified
in~\cite{rfc9110}.

The code is production-shaped and entirely offline. You will build a webhook
receiver that verifies Stripe-style HMAC signatures with a timestamp
tolerance window and an idempotent event store; a delivery worker with
bounded exponential-backoff retries, permanent-4xx short-circuiting, and a
dead-letter log; an SSE producer emitting well-formed
\texttt{text/event-stream} with \texttt{id} and \texttt{retry} fields plus a
dependency-free client parser that resumes with \texttt{Last-Event-ID}; a
WebSocket echo/broadcast demo with protocol-level ping/pong and a clean close
handshake, alongside a pure-stdlib proof of the
\texttt{Sec-WebSocket-Accept} derivation; and an AsyncAPI~3.0 document
describing the Northwind Robotics telemetry contract~\cite{asyncapi}. Every
listing is typed Python~3.11, deterministic, and covered by tests that run
with no network, no credentials, and no clock dependence.

\begin{keyconceptbox}
Polling asks \emph{``has anything happened?''} on the client's schedule; push
delivers \emph{``this happened''} on the server's schedule. The three push
mechanisms are not interchangeable: \textbf{webhooks} push between
\emph{servers} (durable, signed, retried, at-least-once), \textbf{SSE} pushes
a one-way text stream to \emph{browsers and simple clients} over plain HTTP
with automatic reconnection, and \textbf{WebSockets} open a full-duplex
frame-oriented channel for genuinely bidirectional interaction. Choose by
direction of data flow and who the consumer is --- then engineer delivery
semantics, authentication, and backpressure explicitly, because none of the
three gives you those for free.
\end{keyconceptbox}

%% ----------------------------------------------------------------------------
\section{Learning Objectives \& Prerequisites}
\label{sec:ch11-objectives}

\subsection*{Learning Objectives}
After completing this chapter, its lab, and its exercises, you will be able to:

\begin{enumerate}[label=\textbf{LO-\arabic*.}, leftmargin=2.6cm]
  \item Classify any real-time requirement --- dashboards, notifications,
        chat, collaborative editing, server-to-server eventing --- into
        request/response, short polling, long polling, SSE, WebSockets, or
        webhooks, and defend the classification with a latency / fan-out /
        infrastructure trade-off matrix.
  \item Specify a production webhook contract end to end: at-least-once
        delivery semantics, bounded retry schedules with exponential backoff,
        HMAC-SHA256 signature headers in the \texttt{t=...,v1=...} scheme,
        timestamp tolerance windows, endpoint verification handshakes, secret
        rotation with overlapping \texttt{v1} entries, ordering caveats, and
        dead-letter handling.
  \item Implement an idempotent webhook receiver that verifies Stripe-style
        HMAC signatures with a timestamp tolerance window, and a delivery
        worker that distinguishes retryable from permanent failures --- in
        typed Python~3.11 with FastAPI~\cite{fastapidocs}.
  \item Read and write the \texttt{text/event-stream} wire format exactly:
        \texttt{event}/\texttt{data}/\texttt{id}/\texttt{retry} fields,
        multi-line \texttt{data} folding, comment keep-alives; explain
        \texttt{EventSource} reconnection with \texttt{Last-Event-ID}; and
        neutralize buffering proxies with \texttt{X-Accel-Buffering: no}.
  \item Derive \texttt{Sec-WebSocket-Accept} from
        \texttt{Sec-WebSocket-Key} using SHA-1 and the RFC~6455 GUID; parse a
        WebSocket frame header byte by byte (FIN, opcode, mask bit, 7/16/64-bit
        payload lengths); explain why client-to-server masking exists; and use
        ping/pong and close codes correctly.
  \item Decide when WebSockets is the wrong choice --- stateless CRUD,
        cacheable content, hostile intermediaries, horizontally scaled
        stateless services --- and articulate what SSE, long polling, or plain
        REST does better in each case.
  \item Describe an event-driven contract in AsyncAPI~3.0 (channels,
        messages, operations, bindings) and govern event schemas with the
        same rigor as OpenAPI-governed REST contracts~\cite{asyncapi}.
  \item Design consumer-side processing for at-least-once streams:
        idempotent consumers, poison-message quarantine, and the operational
        telemetry that detects retry storms before they become incidents.
\end{enumerate}

\subsection*{Prerequisites}
This chapter assumes you have completed:

\begin{itemize}
  \item \textbf{Chapter 0} --- environment setup: Python 3.11+, a virtual
        environment, and \texttt{pip}/\texttt{pytest} under PowerShell.
  \item \textbf{Chapter 1} --- the HTTP wire protocol: methods, status codes,
        header fields, connection management, and the Upgrade mechanism
        \cite{rfc9110}.
  \item \textbf{Chapter 2} --- consuming APIs from Python, including
        streaming response bodies with \texttt{httpx}~\cite{httpxdocs}.
  \item \textbf{Chapters 3--4} --- REST resource modeling and serialization
        contracts; events are just representations pushed instead of pulled.
  \item \textbf{Chapter 5} --- client-side security fundamentals: TLS,
        bearer tokens~\cite{rfc6749}, and HMAC concepts, which this chapter
        extends to message signing.
  \item \textbf{Chapter 6} --- building APIs with FastAPI: routing, Pydantic
        models, dependency injection, and \texttt{TestClient}
        \cite{fastapidocs}.
  \item \textbf{Chapter 7} --- error modeling with problem details, reused
        here for webhook rejection responses.
\end{itemize}

No network access, cloud account, or API key is required anywhere in this
chapter. Every server binds to \texttt{127.0.0.1}; every fixture is
synthetic, embedded, and deterministic, set in the fictional \emph{Northwind
Robotics} fleet-telemetry universe.

\begin{warningbox}
Webhooks invert the trust relationship you are used to. When you \emph{send}
webhooks you are an HTTP client exposed to arbitrary customer endpoints
(SSRF, slow receivers, retry storms); when you \emph{receive} them you are an
unauthenticated-by-default server that the public internet will probe within
minutes of deployment. Both directions require explicit security
engineering --- signatures, tolerances, timeouts, and bounded retries ---
none of which frameworks provide by default.
\end{warningbox}

%% ----------------------------------------------------------------------------
\section{Deep Technical Mechanics \& Architectural Concepts}
\label{sec:ch11-mechanics}

\subsection{The Interaction-Pattern Taxonomy}
\label{sec:ch11-taxonomy}

Strip away branding and every real-time mechanism reduces to one question:
\emph{who initiates, and how long does the connection live?}

\begin{definition}[Interaction pattern]
An \emph{interaction pattern} is the contract about initiation and connection
lifetime between a consumer and a producer of API data. The four canonical
patterns are: \textbf{request/response} (consumer initiates; one response per
request; connection reusable but logically stateless); \textbf{short polling}
(consumer re-issues request/response on a fixed interval, usually receiving
``nothing new''); \textbf{long polling} (consumer initiates; producer
\emph{holds} the request open until data is available or a deadline passes,
then the consumer immediately re-issues); and \textbf{server push} (producer
initiates delivery over a long-lived or provider-initiated channel: SSE,
WebSocket, or an outbound webhook call).
\end{definition}

Short polling is request/response with a timer, and its economics are brutal
at scale: if a dashboard polls every $5\,\mathrm{s}$ for events that arrive
twice an hour, over $99\%$ of requests are empty work --- yet the load
balancer, TLS terminator, auth middleware, and application worker pay full
price for every one. Long polling converts that waste into connection
occupancy: fewer, longer requests, each of which returns the instant data
arrives. It works through every proxy and firewall on earth because it
\emph{is} plain HTTP, but each notification still costs a full
request/response cycle, and a ``thundering reconnect'' after every publish
can stampede the server.

Server push inverts initiation. SSE keeps one HTTP response open forever and
appends events to it --- strictly server-to-client, text-only, but otherwise
an ordinary HTTP response that enjoys CORS, cookies, and standard auth.
WebSockets perform an HTTP Upgrade and then abandon HTTP semantics entirely
for a bidirectional frame protocol. Webhooks abandon the shared connection
altogether: the \emph{provider} makes a fresh outbound HTTP POST to an
endpoint the \emph{consumer} registered, making webhooks the only pattern
that works server-to-server with no consumer-side listening socket.

Figure~\ref{fig:ch11-patterns} draws the four shapes side by side.

\begin{figure}[htbp]
  \centering
  \begin{tikzpicture}[
      font=\small,
      >=stealth,
      lifeline/.style={draw=packtgray, dashed},
      actor/.style={rectangle, draw=packtgray, fill=lightgray, minimum width=1.7cm, minimum height=0.7cm, align=center},
      flow/.style={->, thick, packtgray},
      pushflow/.style={->, thick, packtorange},
    ]
    % --- Panel 1: request/response ---
    \begin{scope}
      \node[actor] (c1) {Client};
      \node[actor, right=2.2cm of c1] (s1) {Server};
      \draw[lifeline] (c1.south) -- ++(0,-1.5);
      \draw[lifeline] (s1.south) -- ++(0,-1.5);
      \draw[flow] ([yshift=-0.35cm]c1.south) -- node[above, note]{request} ([yshift=-0.35cm]s1.south);
      \draw[flow] ([yshift=-0.85cm]s1.south) -- node[above, note]{response} ([yshift=-0.85cm]c1.south);
      \node[note, below=1.65cm of c1, align=center] {\textbf{Request/response}\\ one shot, client initiates};
    \end{scope}
    % --- Panel 2: short polling ---
    \begin{scope}[xshift=6.2cm]
      \node[actor] (c2) {Client};
      \node[actor, right=2.2cm of c2] (s2) {Server};
      \draw[lifeline] (c2.south) -- ++(0,-1.5);
      \draw[lifeline] (s2.south) -- ++(0,-1.5);
      \draw[flow] ([yshift=-0.25cm]c2.south) -- ([yshift=-0.25cm]s2.south);
      \draw[flow] ([yshift=-0.55cm]s2.south) -- node[above, note]{204 / empty} ([yshift=-0.55cm]c2.south);
      \draw[flow] ([yshift=-0.95cm]c2.south) -- ([yshift=-0.95cm]s2.south);
      \draw[flow] ([yshift=-1.25cm]s2.south) -- node[above, note]{empty again} ([yshift=-1.25cm]c2.south);
      \node[note, below=1.65cm of c2, align=center] {\textbf{Short polling}\\ interval $t$, mostly wasted};
    \end{scope}
    % --- Panel 3: long polling ---
    \begin{scope}[yshift=-4.4cm]
      \node[actor] (c3) {Client};
      \node[actor, right=2.2cm of c3] (s3) {Server};
      \draw[lifeline] (c3.south) -- ++(0,-1.7);
      \draw[lifeline] (s3.south) -- ++(0,-1.7);
      \draw[flow] ([yshift=-0.3cm]c3.south) -- node[above, note]{GET (held open)} ([yshift=-0.3cm]s3.south);
      \draw[decorate, decoration={brace, amplitude=4pt}, packtgray] ([yshift=-0.45cm, xshift=0.1cm]s3.south) -- ([yshift=-1.05cm, xshift=0.1cm]s3.south) node[midway, right=4pt, note]{server waits};
      \draw[pushflow] ([yshift=-1.15cm]s3.south) -- node[above, note]{event arrives} ([yshift=-1.15cm]c3.south);
      \node[note, below=1.85cm of c3, align=center] {\textbf{Long polling}\\ held request, immediate re-issue};
    \end{scope}
    % --- Panel 4: server push ---
    \begin{scope}[xshift=6.2cm, yshift=-4.4cm]
      \node[actor] (c4) {Client};
      \node[actor, right=2.2cm of c4] (s4) {Server};
      \draw[lifeline] (c4.south) -- ++(0,-1.7);
      \draw[lifeline] (s4.south) -- ++(0,-1.7);
      \draw[flow] ([yshift=-0.3cm]c4.south) -- node[above, note]{open channel} ([yshift=-0.3cm]s4.south);
      \draw[pushflow] ([yshift=-0.8cm]s4.south) -- node[above, note]{event 1} ([yshift=-0.8cm]c4.south);
      \draw[pushflow] ([yshift=-1.2cm]s4.south) -- node[above, note]{event 2} ([yshift=-1.2cm]c4.south);
      \node[note, below=1.85cm of c4, align=center] {\textbf{Server push (SSE / WS / webhook)}\\ one open channel, many events};
    \end{scope}
  \end{tikzpicture}
  \caption{The four interaction patterns. Request/response and short polling
  put the client in charge of timing; long polling holds the request until
  the server has something to say; push keeps a channel open (or calls back
  out, for webhooks) so the server initiates delivery.}
  \label{fig:ch11-patterns}
\end{figure}

\begin{table}[htbp]
  \centering
  \small
  \begin{tabularx}{\textwidth}{l X X X X}
    \toprule
    \textbf{Pattern} & \textbf{Typical latency} & \textbf{Fan-out cost} & \textbf{Infrastructure fit} & \textbf{Failure semantics} \\
    \midrule
    Request/response & one round trip & one request per consumer per read & universal; cacheable & clean: status code per call \\
    Short polling & up to interval $t$ & $N$ consumers $\times$ $1/t$ requests, mostly empty & universal; cacheable in principle & clean; but load spikes \\
    Long polling & near-instant per event & one held connection per consumer; reconnect stampede per event & excellent; passes any proxy & held-request timeouts; retry is client logic \\
    SSE & near-instant, continuous & one socket per consumer; cheap frames & good; blocked by buffering proxies; HTTP/1.1 six-per-origin cap & automatic reconnect + \texttt{Last-Event-ID} \\
    WebSocket & near-instant, both directions & one socket per consumer; frame overhead minimal & fair; Upgrade blocked by some proxies; not cacheable & manual reconnect; app-level heartbeat \\
    Webhooks & seconds (queued, retried) & provider makes one outbound POST per subscriber per event & server-to-server only; consumer needs a public endpoint & at-least-once; DLQ after bounded retries \\
    \bottomrule
  \end{tabularx}
  \caption{Pattern selection matrix. ``Latency'' is data-arrival latency seen
  by the consumer; ``fan-out cost'' is the marginal cost of the $N$-th
  consumer.}
  \label{tab:ch11-pattern-matrix}
\end{table}

Two cells of Table~\ref{tab:ch11-pattern-matrix} decide most real
architectures. First: \emph{is the consumer a browser or a server?} Browsers
cannot receive inbound connections, so server-to-server push (webhooks) is
out; browsers speak SSE and WebSocket natively. Second: \emph{is the traffic
one-way or two-way?} A telemetry dashboard is one-way --- SSE wins on
simplicity, reconnection semantics, and proxy friendliness. A collaborative
editor is two-way --- WebSocket is the only honest fit. Polling survives
everywhere as the zero-infrastructure fallback, and Chapter~10's caching
machinery can make short polling of \emph{cacheable} resources nearly free,
which is why ``just poll with \texttt{ETag}'' remains the right answer more
often than fashion suggests.

\subsection{Webhooks: The Reverse-API Pattern}
\label{sec:ch11-webhooks}

\begin{definition}[Webhook]
A \emph{webhook} is a provider-initiated HTTP callback: the consumer
registers an endpoint URL with the provider; when a subscribed event occurs,
the provider serializes the event and POSTs it to that URL. Delivery is
\emph{at-least-once}: the provider retries until it receives a 2xx or
exhausts its retry budget, so duplicates are a contractual certainty, not an
edge case.
\end{definition}

Webhooks look trivial --- ``just POST JSON when something happens'' --- and
that illusion is responsible for a disproportionate share of production
incidents. The provider is now an HTTP client making unsolicited calls to
arbitrary internet endpoints; the consumer is now an unauthenticated server
accepting unsolicited calls from the internet. Every guarantee that plain
REST gets for free must be rebuilt explicitly.

\subsubsection{Delivery semantics and retry schedules}

At-least-once delivery is not a lazy implementation choice; it is the only
semantics compatible with unreliable networks. Exactly-once delivery over a
network partition is provably impossible (the Two Generals problem); the
provider cannot distinguish ``consumer crashed before processing'' from
``consumer processed but the 200 was lost in transit.'' Since the second case
\emph{must} be retried, the first case is retried too --- hence duplicates.
The honest contract is: \emph{at-least-once transport, exactly-once
processing}, where the consumer achieves the second half via idempotency
(Chapter~10).

A production retry schedule is bounded, exponential, and jittered. A typical
shape: attempts at $0$, $1$, $2$, $4$, $8$ minutes, then hourly for up to 24
or 72 hours, then dead-letter. Jitter (random $\pm 20\%$) prevents
synchronized retry storms when a popular consumer recovers from an outage and
ten thousand queued deliveries fire simultaneously. Crucially, the schedule
distinguishes \emph{retryable} failures --- connection errors, timeouts, 408,
409, 425, 429, and 5xx --- from \emph{permanent} ones: a 400 or 422 means the
payload itself is rejected and retrying it ten thousand times is an
own-goal DDoS (Section~\ref{sec:ch11-lab} tells that war story).

\subsubsection{Signature verification: the \texttt{t=...,v1=...} scheme}

A webhook endpoint is a public URL that accepts POSTs from anyone who finds
it. Without authentication, whoever guesses the URL can inject fabricated
``payment succeeded'' or ``robot docked'' events. Bearer tokens in a
configured header help, but the industry standard --- popularized by Stripe
and now near-universal --- is symmetric \textbf{HMAC-SHA256 message signing}:

\begin{equation}
  \sigma \;=\; \operatorname{HMAC\text{-}SHA256}\bigl(k,\; t \,\|\, \texttt{"."} \,\|\, b\bigr)
  \label{eq:ch11-hmac}
\end{equation}

where $k$ is the per-endpoint signing secret, $t$ is the Unix timestamp at
signing time, and $b$ is the \emph{raw, unparsed} request body. The provider
sends the timestamp and signature together in one header:

\begin{minted}{text}
POST /webhooks/telemetry HTTP/1.1
Host: ops.customer.example
Content-Type: application/json
X-Northwind-Signature: t=1700000000,v1=5257a869e7ecebeda32affa62cdca3fa51cad7e77a0e56ff536d0ce8e108d8bd
                        \_________/ \____________________________________________________/
                        timestamp   hex HMAC-SHA256 over "1700000000.<raw-body>"
                        (replay     (proves possession of the shared secret k;
                         anchor)     multiple v1= entries allowed during rotation)
\end{minted}

The receiver recomputes Equation~\eqref{eq:ch11-hmac} over the raw body and
the header's own $t$, then compares with constant-time semantics
(\texttt{hmac.compare\_digest}), because a naive \texttt{==} on hex strings
leaks prefix-match timing that can be iterated into a forgery oracle. Three
independent checks, three independent failure modes:

\begin{enumerate}
  \item \textbf{Freshness:} $|now - t| \leq$ tolerance (typically 300\,s).
        This is the replay-protection window: a captured request replayed
        later carries a stale $t$ and is rejected regardless of signature
        validity. The tolerance exists only to absorb clock skew and delivery
        latency; keep it tight.
  \item \textbf{Authenticity:} at least one \texttt{v1} candidate matches the
        recomputed HMAC. Multiple \texttt{v1} values in one header are the
        secret-rotation mechanism: during rotation the provider signs with
        both the old and new secret, and receivers configured with either
        accept.
  \item \textbf{Integrity:} the HMAC covers the timestamp and body together,
        so neither can be altered without detection. Signing the raw bytes
        --- not a re-serialized object --- matters: JSON re-serialization is
        not byte-stable (key order, whitespace, Unicode escapes), so verify
        \emph{before} parsing.
\end{enumerate}

\subsubsection{Endpoint verification, rotation, ordering, and the DLQ}

Before a provider sends real events, it must prove the registrant actually
controls the endpoint --- otherwise webhooks become an amplification vector
for attacking third parties. The standard \emph{verification handshake}: at
registration time the provider POSTs a challenge (a random token) and
requires the endpoint to echo it back, or requires the registrant to complete
an out-of-band confirmation. Slack-style \texttt{url\_verification}
challenges and WebSub-style hub-challenge GETs are the canonical shapes.

Secret rotation follows the overlap pattern above: issue $k'$, sign with both
$k$ and $k'$ for one tolerance window, then retire $k$. Providers must also
support consumer-triggered rotation with immediate effect --- when a consumer
suspects a leak, ``rotate within 24 hours'' is not an answer.

\textbf{Ordering is not part of the contract.} Retries, parallel delivery
workers, and per-endpoint queues conspire to deliver \texttt{created} after
\texttt{updated}, or to deliver the same event twice with a third event in
between. Consumers must either be order-insensitive (each event carries full
current state --- the \emph{state-transfer} style) or must buffer and reorder
by a monotonic sequence number or \texttt{occurred\_at} timestamp the
provider includes in every payload. Northwind Robotics events carry both an
\texttt{id} and an \texttt{occurred\_at} for exactly this reason.

Finally, the \textbf{dead-letter queue}: when the retry budget is exhausted,
the event and its full attempt history must land somewhere durable and
alerted, with a replay tool. A webhook provider without a DLQ silently loses
customer data on the first extended consumer outage; a DLQ without alerting
is a write-only graveyard.

\begin{keyconceptbox}
\textbf{Provider checklist:} verify endpoint ownership before first delivery;
sign every delivery (HMAC-SHA256 over timestamp + raw body); retry only
retryable failures, with bounded exponential backoff plus jitter; honor
consumer-driven secret rotation immediately; include a unique event id and
timestamp in every payload; dead-letter with alerting and a replay tool;
document that ordering is best-effort.
\textbf{Consumer checklist:} verify freshness \emph{then} signature, in
constant time, over raw bytes, before parsing; respond 2xx only after the
event is durably recorded (or 2xx-fast with a queue in front); make
processing idempotent on event id; never log raw payloads or secrets;
return 4xx for malformed requests so the provider stops retrying; monitor
duplicate rates --- a spike means a retry storm is in progress.
\end{keyconceptbox}

Figure~\ref{fig:ch11-webhook-fsm} captures the provider-side delivery
lifecycle as a state machine; the code in
Section~\ref{sec:ch11-code-webhooks} implements it exactly.

\begin{figure}[htbp]
  \centering
  \begin{tikzpicture}[
      font=\small,
      >=stealth,
      state/.style={rectangle, rounded corners, draw=packtgray, fill=blue!5, minimum width=2.6cm, minimum height=0.85cm, align=center},
      terminal/.style={rectangle, rounded corners, draw=packtgray, fill=green!10, minimum width=2.6cm, minimum height=0.85cm, align=center},
      dead/.style={rectangle, rounded corners, draw=packtgray, fill=yellow!15, minimum width=2.6cm, minimum height=0.85cm, align=center},
      edge/.style={->, thick, packtgray},
      every label/.style={note},
    ]
    \node[state] (queued) {QUEUED};
    \node[state, right=2.4cm of queued] (attempt) {ATTEMPTING};
    \node[terminal, right=2.4cm of attempt] (delivered) {DELIVERED};
    \node[state, below=1.5cm of attempt] (backoff) {BACKOFF-WAIT};
    \node[dead, left=2.4cm of backoff] (dlq) {DEAD-LETTERED};

    \draw[edge] (queued) -- node[above, note]{worker picks up} (attempt);
    \draw[edge] (attempt) -- node[above, note]{2xx} (delivered);
    \draw[edge] (attempt) -- node[right, note, align=left]{retryable failure\\ (5xx / 429 / timeout)\\ and attempts $<$ max} (backoff);
    \draw[edge] (backoff) -- node[above, note]{delay $2^{n}$ expires} (attempt);
    \draw[edge] (attempt) -| node[pos=0.25, below, note, align=center]{permanent 4xx\\ or attempts $=$ max} (dlq);
    \draw[edge] (dlq) -- node[left, note, align=right]{operator replay\\ (re-enqueue)} (queued);
  \end{tikzpicture}
  \caption{Webhook delivery state machine. Every terminal failure is either a
  permanent rejection (do not retry a payload the consumer has told you is
  invalid) or budget exhaustion into the dead-letter queue, from which only
  an explicit operator action restores the event.}
  \label{fig:ch11-webhook-fsm}
\end{figure}

\subsection{Server-Sent Events: \texttt{text/event-stream} on the Wire}
\label{sec:ch11-sse}

Server-Sent Events are the most underrated mechanism in this chapter: a
completely standard HTTP response whose body simply never ends, carrying a
tiny text protocol that browsers parse natively via the \texttt{EventSource}
API. There is no upgrade, no new framing layer, no new attack surface --- the
stream is cache-control-compatible, CORS-compatible, cookie-authenticated,
and readable with \texttt{curl}.

\begin{definition}[Server-Sent Events]
\emph{Server-Sent Events} (SSE) is a W3C/WHATWG protocol in which a server
responds to a GET with \texttt{Content-Type: text/event-stream} and then
writes a sequence of UTF-8 \emph{event blocks} separated by blank lines. Each
block consists of \texttt{field: value} lines; the defined fields are
\texttt{event} (named type, default \texttt{message}), \texttt{data} (the
payload; multiple \texttt{data} lines are joined with \texttt{\textbackslash
n}), \texttt{id} (the last-event-id cursor used for reconnection), and
\texttt{retry} (reconnection delay in milliseconds). Lines beginning with a
colon are comments and serve as keep-alives.
\end{definition}

The wire format, exactly as bytes on the connection:

\begin{minted}{text}
HTTP/1.1 200 OK
Content-Type: text/event-stream
Cache-Control: no-cache
X-Accel-Buffering: no

: connected to northwind telemetry (comment line = keep-alive; proxies flush on bytes)

event: telemetry
data: {"robot": "NWR-07",
data: "battery": 0.82}
id: 128
retry: 2000

event: alert
data: {"robot": "NWR-07", "alert": "dock requested"}
id: 129

: keep-alive (sent every ~15 s so idle proxies do not reap the connection)
\end{minted}

The grammar rules that bite in production: field name and value are separated
by a colon plus one \emph{optional} leading space; a blank line
\textbf{dispatches} the buffered event; a block with no \texttt{data} field
is discarded; a \texttt{data} value split across lines is rejoined with
newlines in order; \texttt{id} must not contain newlines and \emph{persists}
--- the client's last-event-id is sticky until replaced; \texttt{retry} is
parsed as ASCII digits only, everything else is ignored.

\subsubsection{Reconnection and \texttt{Last-Event-ID}}

SSE's killer feature is that reconnection is \emph{protocol}, not
application code. When the TCP connection drops, the browser's
\texttt{EventSource} waits the most recent \texttt{retry} interval (default
three seconds), re-issues the GET, and automatically sets the
\texttt{Last-Event-ID} request header to the id of the last dispatched event.
A correct server reads that header and resumes the stream \emph{after} it ---
from a ring buffer, an event log, or a broker offset. This gives SSE
something WebSocket does not have: a standardized, zero-code resumability
story. The contract has one sharp edge: the server decides how much history
to retain, so ``resume'' is bounded by the retention window; a client that
reconnects after the window has scrolled past must receive a signal (a
\texttt{control} event, or an HTTP 204 telling \texttt{EventSource} to stop)
and fall back to a full re-fetch.

\begin{examplebox}
The browser side is deliberately unremarkable --- that is the point:
\begin{minted}{text}
const source = new EventSource("/stream");          // cookies flow automatically
source.addEventListener("telemetry", (e) => render(JSON.parse(e.data)));
source.addEventListener("alert",     (e) => escalate(JSON.parse(e.data)));
source.onerror = () => { /* browser retries automatically with Last-Event-ID */ };
// e.lastEventId is the id: field; server-set "retry: 2000" controls the backoff
\end{minted}
No reconnection logic, no heartbeat logic, no framing logic. Everything that
WebSocket leaves to you, the SSE user agent already does.
\end{examplebox}

\subsubsection{Proxy buffering and HTTP/2}

Two infrastructure facts kill more SSE deployments than any protocol bug.
First, \textbf{buffering intermediaries}: many reverse proxies (nginx in
default proxy mode, some CDNs, some corporate egress proxies) buffer response
bodies before forwarding --- which turns your real-time stream into a
not-real-time dump delivered when the buffer fills or the connection closes.
The defenses: set \texttt{X-Accel-Buffering: no} (honored by nginx and
Envoy), send \texttt{Cache-Control: no-cache}, flush after every event, and
emit comment keep-alives every 10--15 seconds so idle-timeout middleboxes see
traffic. Second, \textbf{the HTTP/1.1 connection cap}: browsers allow only
six concurrent connections per origin, and each \texttt{EventSource} holds
one for its entire lifetime --- two open dashboard tabs can starve every
other request to that origin. Over HTTP/2 the cap effectively disappears
(streams are multiplexed over one connection), so SSE \emph{wants} HTTP/2 ---
with the caveat that a single TCP connection's head-of-line blocking then
applies to everything multiplexed on it.

\subsection{WebSockets: RFC~6455 Precisely}
\label{sec:ch11-websockets}

WebSocket is not ``HTTP but faster.'' It is a distinct, frame-oriented,
bidirectional protocol that borrows HTTP only for its opening handshake ---
deliberately, so that the same port 80/443 infrastructure, TLS termination,
and Upgrade machinery~\cite{rfc9110} can be reused.

\subsubsection{The opening handshake}

The client sends an ordinary HTTP/1.1 GET with \texttt{Upgrade: websocket},
\texttt{Connection: Upgrade}, a protocol version (\texttt{Sec-WebSocket-
Version: 13}), and a randomly generated 16-byte nonce, Base64-encoded, in
\texttt{Sec-WebSocket-Key}. The server proves it understood the WebSocket
upgrade --- and was not merely an HTTP server echoing headers --- by
computing:

\begin{equation}
  \texttt{Sec-WebSocket-Accept} \;=\;
  \operatorname{Base64}\Bigl(\operatorname{SHA\text{-}1}\bigl(\texttt{Sec-WebSocket-Key} \,\|\, \texttt{"258EAFA5-E914-47DA-95CA-C5AB0DC85B11"}\bigr)\Bigr)
  \label{eq:ch11-ws-accept}
\end{equation}

The magic GUID is fixed by RFC~6455 section~1.3; it guarantees the
concatenated string can never be a plausible HTTP artifact, so a
non-WebSocket server cannot accidentally produce a valid accept value. The
worked example from the RFC itself --- which Listing~11.6
reproduces and asserts in pure stdlib Python --- is:

\begin{minted}{text}
Client:  Sec-WebSocket-Key: dGhlIHNhbXBsZSBub25jZQ==
Server:  SHA-1("dGhlIHNhbXBsZSBub25jZQ==258EAFA5-E914-47DA-95CA-C5AB0DC85B11")
       = 0xb3 0x7a 0x4f 0x2c 0xc0 0x62 0x4f 0x16 0x90 0xf6
         0x46 0x06 0xcf 0x38 0x59 0x45 0xb2 0xbe 0xc4 0xea
       -> Base64 -> Sec-WebSocket-Accept: s3pPLMBiTxaQ9kYGzzhZRbK+xOo=
\end{minted}

The server answers \texttt{101 Switching Protocols}, and from that byte
onward the connection speaks frames, not HTTP. Optional handshake headers
negotiate \emph{subprotocols} (\texttt{Sec-WebSocket-Protocol:
mqtt, chat.v2}) and extensions (\texttt{Sec-WebSocket-Extensions:
permessage-deflate} for compression); the server picks at most one
subprotocol, and a client that required one must fail the connection if the
server's pick is absent.

\subsubsection{Frame format}

Every post-handshake message travels in frames with a 2--14 byte header:

\begin{minted}{text}
 byte 0                     byte 1                     extended length
+-+-+-+-+-+-+-+-----------+-+-----------+-----------------------------+
|F|R|R|R|opcode|M| payload len (7 bits) |  0 / 16 / 64 extra bits     |
|I|S|S|S|       |A|                       (126 -> 16-bit,             |
|N|V|V|V|       |S|                        127 -> 64-bit big-endian)  |
+-+-+-+-+-+-+-+-----------+-+-----------+-----------------------------+
| masking key (32 bits, present iff MASK=1; client->server MUST mask) |
+---------------------------------------------------------------------+
| payload (XORed with masking key when masked)                        |
+---------------------------------------------------------------------+

opcodes: 0x0 continuation  0x1 text (UTF-8)  0x2 binary
         0x8 close         0x9 ping          0xA pong    (0x3-0x7, 0xB-0xF reserved)
\end{minted}

The header fields: \textbf{FIN} marks the final fragment of a message
(messages may be fragmented into continuation frames, so receivers must
reassemble); \textbf{RSV1--3} are zero unless an extension (e.g.\
permessage-deflate) claims them; the 4-bit \textbf{opcode} selects data
versus control; the \textbf{MASK} bit announces whether a masking key
follows; and the payload length is \emph{self-describing} --- values 0--125
are literal, 126 means ``read the next 2 bytes as a big-endian uint16,'' and
127 means ``read the next 8 bytes as a big-endian uint63'' (the high bit must
be zero). Control frames (close, ping, pong) must have FIN set and payloads
$\leq 125$ bytes, and may be interleaved \emph{between} fragments of a data
message.

\textbf{Why masking exists} is the most-asked WebSocket interview question,
and the answer is not confidentiality. Client-to-server frames are XOR-masked
with a fresh random 32-bit key per frame to defend \emph{intermediaries}, not
endpoints: researchers showed that a hostile script in a browser could
otherwise craft bytes that, to a misbehaving transparent proxy, look like a
smuggled HTTP request --- poisoning shared caches (the 2010 ``cache
poisoning'' attack that motivated the mask). Because the key is chosen by the
sender's stack after the application supplied the bytes, the attacker cannot
predict what the masked output will look like on the wire. Masking costs a
few cycles; server-to-client frames skip it.

\textbf{Ping/pong} (opcodes 0x9/0xA) is protocol-level keep-alive and
liveness detection: either peer may ping; the other must pong ``as soon as
practical'' with identical payload. TCP itself will happily keep a dead
connection ``open'' for hours, so applications that care about zombie
connections ping every 20--30 seconds and close after a missed pong.
\textbf{Close} (0x8) is a handshake, not a hang-up: the initiator sends a
close frame with a 2-byte status code and optional UTF-8 reason; the peer
echoes a close frame; then the TCP connection drops. Key codes: 1000 normal,
1001 going away (server restart / page navigation), 1002 protocol error,
1003 unsupported data, 1006 is \emph{reserved} for abnormal closure and must
never appear on the wire, 1007 invalid payload, 1008 policy violation, 1009
message too big, 1011 internal server error.

\subsubsection{When WebSockets is the wrong choice}

WebSocket's power is exactly its cost: it abandons HTTP semantics, so nothing
in the HTTP toolchain applies --- no caching, no \texttt{ETag}, no
idempotency vocabulary, no status codes, no content negotiation, no
meaningful gateway routing. Reach for something else when:

\begin{itemize}
  \item \textbf{The work is stateless CRUD.} Create/read/update/delete with
        request/response semantics wants REST: you get caching, retries,
        tracing, and firewall friendliness for free~\cite{rfc9110}.
  \item \textbf{The data is cacheable.} A WebSocket frame cannot be cached by
        a CDN; an \texttt{ETag}-polled resource can be nearly free at scale.
  \item \textbf{The traffic is one-way.} SSE is simpler, resumable, and
        proxy-friendlier; WebSocket buys you a send channel you will never
        use and bills you in connection-management code.
  \item \textbf{Hostile intermediaries are in the path.} Some corporate
        proxies and older load balancers mangle or refuse \texttt{Upgrade}.
        SSE and long polling degrade gracefully; WebSocket just fails.
  \item \textbf{You need horizontal statelessness.} Each WebSocket pins a
        client to a server process; broadcast across a fleet requires a
        backplane (broker fan-out, Redis pub/sub). SSE has the same pinning
        but simpler recovery, because reconnection with
        \texttt{Last-Event-ID} is built in.
\end{itemize}

Figure~\ref{fig:ch11-topology} contrasts the two push topologies directly.

\begin{figure}[htbp]
  \centering
  \begin{tikzpicture}[
      font=\small,
      >=stealth,
      box/.style={rectangle, draw=packtgray, fill=lightgray, minimum width=2.0cm, minimum height=0.8cm, align=center},
      srv/.style={rectangle, draw=packtgray, fill=blue!5, minimum width=2.2cm, minimum height=0.9cm, align=center},
      conn/.style={<->, thick, packtorange},
      uni/.style={->, thick, packtorange},
      lbl/.style={note, align=center},
    ]
    % SSE panel
    \node[srv] (sse) {SSE server\\ (plain HTTP)};
    \node[box, above left=1.2cm and 1.6cm of sse] (b1) {Browser tab};
    \node[box, above=1.2cm of sse] (b2) {Browser tab};
    \node[box, above right=1.2cm and 1.6cm of sse] (b3) {\texttt{httpx} client};
    \draw[uni] (sse) -- node[lbl, pos=0.55, above left=-2pt]{\texttt{GET} then \texttt{text/event-stream}\\ one-way; auto-reconnect + Last-Event-ID} (b1);
    \draw[uni] (sse) -- (b2);
    \draw[uni] (sse) -- (b3);
    \node[lbl, below=0.35cm of sse] {\textbf{SSE topology:} $N$ long-lived HTTP responses, server $\to$ client only};

    % WebSocket panel
    \node[srv, right=6.8cm of sse] (ws) {WebSocket server\\ (post-Upgrade)};
    \node[box, above left=1.2cm and 1.6cm of ws] (w1) {Browser tab};
    \node[box, above=1.2cm of ws] (w2) {Browser tab};
    \node[box, above right=1.2cm and 1.6cm of ws] (w3) {Robot console};
    \draw[conn] (ws) -- node[lbl, pos=0.55, above left=-2pt]{full-duplex frames\\ text/binary, ping/pong} (w1);
    \draw[conn] (ws) -- (w2);
    \draw[conn] (ws) -- (w3);
    \node[lbl, below=0.35cm of ws] {\textbf{WebSocket topology:} $N$ upgraded channels, both directions};
  \end{tikzpicture}
  \caption{SSE versus WebSocket topology. Both pin one socket per client;
  SSE's channel is a plain HTTP response with standardized resume, while
  WebSocket's is a symmetric frame channel requiring application-level
  reconnection.}
  \label{fig:ch11-topology}
\end{figure}

\subsection{AsyncAPI and the Governance of Event Contracts}
\label{sec:ch11-asyncapi}

REST APIs got their contract discipline from OpenAPI; event-driven APIs get
theirs from \textbf{AsyncAPI}~\cite{asyncapi}. The conceptual mapping is
direct: where OpenAPI has paths and operations, AsyncAPI has \emph{channels}
(addressable destinations such as a Kafka topic
\texttt{northwind.robots.\{robot\_id\}.telemetry}) and \emph{operations}
(send/receive over a channel); where OpenAPI has request/response schemas,
AsyncAPI has \emph{messages} with payload schemas; \emph{bindings} attach
protocol-specific configuration (Kafka partitions, MQTT QoS, WebSocket
subprotocols) without polluting the logical contract. Listing~11.8
shows the Northwind Robotics telemetry contract in AsyncAPI~3.0.

Governance is the point, not the YAML. Event schemas evolve under exactly the
same compatibility algebra as REST representations (Chapter~9): additive
changes are safe, removals and renamings break consumers you have never met
--- and in event-driven systems you \emph{cannot} meet them, because the
whole point of the broker is that producers do not know their consumers.
Discipline therefore means: a schema registry or repository as the single
source of truth, CI gates that reject breaking schema diffs, explicit
compatibility modes (backward for consumers-first rollout), and a deprecation
workflow for event types exactly as for endpoints. The consumer-side
corollaries mirror the webhook contract from
Section~\ref{sec:ch11-webhooks}: assume \textbf{at-least-once} delivery and
make every consumer idempotent on the event id; treat \textbf{poison
messages} --- events that crash or reject deterministically --- as a routing
problem (quarantine to a dead-letter topic after $n$ attempts) rather than an
infinite retry; and never trust the payload beyond the schema, because a
broker is a trust boundary like any other~\cite{kleppmann2017ddia,
newman2021microservices}.

%% ----------------------------------------------------------------------------
\section{Production-Ready Code Implementation}
\label{sec:ch11-code}

All listings run on Python~3.11+ (verified on 3.13) with the dependencies in
Section~\ref{sec:ch11-deps}. Everything binds to \texttt{127.0.0.1} or uses
in-process ASGI transport: no network, no credentials, no wall-clock
dependence. PowerShell setup:

\begin{minted}{text}
python -m venv .venv
.\.venv\Scripts\Activate.ps1
pip install -r requirements.txt
pytest -q
\end{minted}

\subsection{Webhooks: Signed Receiver and Bounded-Retry Delivery Worker}
\label{sec:ch11-code-webhooks}

Listing~11.1 implements the consumer side of
Section~\ref{sec:ch11-webhooks}: freshness check, constant-time signature
comparison over the raw body, and an idempotency store that turns
at-least-once delivery into exactly-once processing. The clock is injectable
so tests are deterministic.

\begin{minted}{python}
"""Webhook receiver with Stripe-style HMAC-SHA256 signature verification."""
from __future__ import annotations

import hashlib
import hmac
import json
import time
from collections.abc import Callable
from dataclasses import dataclass, field
from typing import Any

from fastapi import FastAPI, Request
from fastapi.responses import JSONResponse

SIGNATURE_HEADER = "x-northwind-signature"
TOLERANCE_SECONDS = 300


def compute_signature(secret: str, timestamp: int, body: bytes) -> str:
    """Return the hex HMAC-SHA256 of "<timestamp>.<body>"."""
    signed_payload = str(timestamp).encode("utf-8") + b"." + body
    digest = hmac.new(secret.encode("utf-8"), signed_payload, hashlib.sha256)
    return digest.hexdigest()


def parse_signature_header(header: str) -> tuple[int, list[str]]:
    """Parse "t=...,v1=..." into (timestamp, [signatures]); ValueError if malformed."""
    timestamp: int | None = None
    signatures: list[str] = []
    for part in header.split(","):
        key, _, value = part.partition("=")
        if key.strip() == "t":
            try:
                timestamp = int(value)
            except ValueError as exc:
                raise ValueError("non-integer timestamp") from exc
        elif key.strip() == "v1":
            signatures.append(value.strip())
    if timestamp is None or not signatures:
        raise ValueError("malformed signature header")
    return timestamp, signatures


@dataclass
class IdempotencyStore:
    """In-memory set of processed event ids (swap for Redis/Postgres in prod)."""

    processed: set[str] = field(default_factory=set)

    def is_duplicate(self, event_id: str) -> bool:
        return event_id in self.processed

    def mark_processed(self, event_id: str) -> None:
        self.processed.add(event_id)


def _problem(status: int, title: str, detail: str) -> JSONResponse:
    return JSONResponse(
        status_code=status,
        media_type="application/problem+json",
        content={"type": "about:blank", "title": title,
                 "status": status, "detail": detail},
    )


def create_app(
    secret: str,
    store: IdempotencyStore,
    now: Callable[[], float] = time.time,
    tolerance_seconds: int = TOLERANCE_SECONDS,
) -> FastAPI:
    app = FastAPI(title="Northwind Robotics Webhook Receiver")

    @app.post("/webhooks/telemetry")
    async def receive_telemetry(request: Request) -> JSONResponse:
        body = await request.body()
        raw_header = request.headers.get(SIGNATURE_HEADER, "")
        try:
            timestamp, signatures = parse_signature_header(raw_header)
        except ValueError:
            return _problem(400, "Malformed Signature",
                            "Signature header is missing or malformed.")

        # Replay protection: refuse events signed too far from our clock.
        if abs(now() - timestamp) > tolerance_seconds:
            return _problem(400, "Stale Timestamp",
                            "Signature timestamp outside tolerance window.")

        expected = compute_signature(secret, timestamp, body)
        if not any(hmac.compare_digest(expected, c) for c in signatures):
            return _problem(400, "Invalid Signature",
                            "No v1 signature matches the payload.")

        try:
            event: dict[str, Any] = json.loads(body)
            event_id = str(event["id"])
        except (json.JSONDecodeError, KeyError, TypeError):
            return _problem(400, "Malformed Event",
                            "Body is not a JSON event with an id field.")

        # Idempotent receiver: at-least-once delivery means duplicates WILL arrive.
        if store.is_duplicate(event_id):
            return JSONResponse(status_code=200,
                                content={"status": "duplicate", "id": event_id})
        store.mark_processed(event_id)
        return JSONResponse(status_code=200,
                            content={"status": "processed", "id": event_id})

    return app
\end{minted}

Listing~11.2 is the provider side: the state machine of
Figure~\ref{fig:ch11-webhook-fsm} in code. Transport, clock, and sleeper are
injected; backoff delays are pure functions of the attempt number, so the
test suite observes the schedule without sleeping.

\begin{minted}{python}
"""Webhook delivery worker: bounded retries, backoff, and a dead-letter log."""
from __future__ import annotations

import json
import time
from collections.abc import Callable
from dataclasses import dataclass
from enum import Enum
from typing import Any

from webhook_receiver import compute_signature

RETRYABLE_STATUSES = frozenset({408, 409, 425, 429, 500, 502, 503, 504})


class DeliveryOutcome(str, Enum):
    DELIVERED = "delivered"
    DEAD_LETTERED = "dead_lettered"
    REJECTED_PERMANENT = "rejected_permanent"


@dataclass(frozen=True)
class WebhookEvent:
    event_id: str
    event_type: str
    occurred_at: int
    payload: dict[str, Any]

    def body_bytes(self) -> bytes:
        doc = {"id": self.event_id, "type": self.event_type,
               "occurred_at": self.occurred_at, "data": self.payload}
        return json.dumps(doc, sort_keys=True).encode("utf-8")


@dataclass(frozen=True)
class Attempt:
    number: int
    status: int | None  # None => transport raised (connection refused, DNS, ...)
    retryable: bool


@dataclass(frozen=True)
class DeadLetter:
    event_id: str
    endpoint: str
    attempts: tuple[Attempt, ...]
    reason: str


Transport = Callable[[str, dict[str, str], bytes], int]


class WebhookDeliverer:
    def __init__(
        self,
        secret: str,
        transport: Transport,
        *,
        max_attempts: int = 5,
        base_delay_seconds: float = 1.0,
        now: Callable[[], float] = time.time,
        sleeper: Callable[[float], None] = time.sleep,
    ) -> None:
        if max_attempts < 1:
            raise ValueError("max_attempts must be >= 1")
        self._secret = secret
        self._transport = transport
        self._max_attempts = max_attempts
        self._base_delay = base_delay_seconds
        self._now = now
        self._sleeper = sleeper
        self.dead_letters: list[DeadLetter] = []

    def backoff_delay(self, attempt_number: int) -> float:
        """Exponential backoff with a deterministic cap: 1, 2, 4, 8, ... capped at 60s."""
        return min(self._base_delay * (2 ** (attempt_number - 1)), 60.0)

    def deliver(self, endpoint: str, event: WebhookEvent) -> DeliveryOutcome:
        body = event.body_bytes()
        attempts: list[Attempt] = []
        for number in range(1, self._max_attempts + 1):
            timestamp = int(self._now())
            headers = {
                "content-type": "application/json",
                "x-northwind-signature": (
                    f"t={timestamp},v1={compute_signature(self._secret, timestamp, body)}"
                ),
            }
            try:
                status = self._transport(endpoint, headers, body)
                retryable = status in RETRYABLE_STATUSES
            except OSError:
                status, retryable = None, True
            attempts.append(Attempt(number=number, status=status, retryable=retryable))

            if status is not None and 200 <= status < 300:
                return DeliveryOutcome.DELIVERED
            if status is not None and not retryable:
                # 4xx (other than 408/409/425/429) means "stop calling me";
                # retrying forever is how webhook loops turn into a DDoS.
                self._dead_letter(event, endpoint, attempts, f"permanent {status}")
                return DeliveryOutcome.REJECTED_PERMANENT
            if number < self._max_attempts:
                self._sleeper(self.backoff_delay(number))

        self._dead_letter(event, endpoint, attempts, "max attempts exhausted")
        return DeliveryOutcome.DEAD_LETTERED

    def _dead_letter(
        self, event: WebhookEvent, endpoint: str, attempts: list[Attempt], reason: str
    ) -> None:
        self.dead_letters.append(
            DeadLetter(event.event_id, endpoint, tuple(attempts), reason)
        )
\end{minted}

The test suite (Listing~11.3) covers the attacks and failure modes that
matter: tampered bodies, stale timestamps, wrong secrets, replayed
deliveries, retry-then-succeed, budget exhaustion into the DLQ, and the
permanent-4xx short-circuit.

\begin{minted}{python}
"""Offline, deterministic tests for the webhook receiver and delivery worker."""
from __future__ import annotations

import json

import pytest
from fastapi.testclient import TestClient

from webhook_delivery import DeliveryOutcome, WebhookDeliverer, WebhookEvent
from webhook_receiver import (
    IdempotencyStore, compute_signature, create_app, parse_signature_header,
)

SECRET = "whsec_northwind_test_only"
NOW = 1_700_000_000  # fixed clock: no wall-time dependence


def signed_headers(secret: str, timestamp: int, body: bytes) -> dict[str, str]:
    sig = compute_signature(secret, timestamp, body)
    return {"x-northwind-signature": f"t={timestamp},v1={sig}"}


def event_body(event_id: str = "evt_1001") -> bytes:
    return json.dumps({"id": event_id, "type": "robot.telemetry",
                       "data": {"robot": "NWR-07", "battery": 0.82}}).encode()


@pytest.fixture()
def client() -> TestClient:
    app = create_app(SECRET, IdempotencyStore(), now=lambda: float(NOW))
    return TestClient(app)


def test_valid_signature_is_processed(client: TestClient) -> None:
    body = event_body()
    resp = client.post("/webhooks/telemetry", content=body,
                       headers=signed_headers(SECRET, NOW, body))
    assert resp.status_code == 200
    assert resp.json() == {"status": "processed", "id": "evt_1001"}


def test_tampered_body_is_rejected(client: TestClient) -> None:
    body = event_body()
    headers = signed_headers(SECRET, NOW, body)
    tampered = body.replace(b"0.82", b"0.01")  # attacker flips the battery level
    resp = client.post("/webhooks/telemetry", content=tampered, headers=headers)
    assert resp.status_code == 400
    assert resp.json()["title"] == "Invalid Signature"


def test_stale_timestamp_is_rejected(client: TestClient) -> None:
    body = event_body()
    old = NOW - 3600  # signed an hour ago: outside the 300 s tolerance
    resp = client.post("/webhooks/telemetry", content=body,
                       headers=signed_headers(SECRET, old, body))
    assert resp.status_code == 400
    assert resp.json()["title"] == "Stale Timestamp"


def test_wrong_secret_is_rejected(client: TestClient) -> None:
    body = event_body()
    headers = signed_headers("whsec_attacker_guess", NOW, body)
    resp = client.post("/webhooks/telemetry", content=body, headers=headers)
    assert resp.status_code == 400


def test_replayed_delivery_is_idempotent(client: TestClient) -> None:
    body = event_body()
    headers = signed_headers(SECRET, NOW, body)
    first = client.post("/webhooks/telemetry", content=body, headers=headers)
    replay = client.post("/webhooks/telemetry", content=body, headers=headers)
    assert first.json()["status"] == "processed"
    assert replay.status_code == 200  # ack again, but do not reprocess
    assert replay.json()["status"] == "duplicate"


def test_header_parsing_supports_rotation_overlap() -> None:
    ts, sigs = parse_signature_header("t=1700000000,v1=aaa,v1=bbb")
    assert ts == NOW
    assert sigs == ["aaa", "bbb"]


def make_event() -> WebhookEvent:
    return WebhookEvent("evt_1001", "robot.telemetry", NOW, {"robot": "NWR-07"})


class FakeTransport:
    """Scripted statuses; records every call. OSError simulates connection failure."""

    def __init__(self, script: list[int | type[OSError]]) -> None:
        self.script = script
        self.calls: list[tuple[str, dict[str, str], bytes]] = []

    def __call__(self, url: str, headers: dict[str, str], body: bytes) -> int:
        self.calls.append((url, headers, body))
        outcome = self.script.pop(0)
        if isinstance(outcome, type):
            raise outcome("connection refused")
        return outcome


def make_deliverer(transport: FakeTransport) -> tuple[WebhookDeliverer, list[float]]:
    sleeps: list[float] = []
    deliverer = WebhookDeliverer(
        SECRET, transport, max_attempts=5, base_delay_seconds=1.0,
        now=lambda: float(NOW), sleeper=sleeps.append,
    )
    return deliverer, sleeps


def test_delivery_succeeds_first_try() -> None:
    transport = FakeTransport([200])
    deliverer, sleeps = make_deliverer(transport)
    outcome = deliverer.deliver("https://ops.northwind-robotics.example/hook", make_event())
    assert outcome is DeliveryOutcome.DELIVERED
    assert len(transport.calls) == 1 and sleeps == [] and deliverer.dead_letters == []


def test_delivery_retries_5xx_with_backoff_then_succeeds() -> None:
    transport = FakeTransport([500, 503, 200])
    deliverer, sleeps = make_deliverer(transport)
    outcome = deliverer.deliver("https://ops.northwind-robotics.example/hook", make_event())
    assert outcome is DeliveryOutcome.DELIVERED
    assert len(transport.calls) == 3
    assert sleeps == [1.0, 2.0]  # deterministic exponential backoff, no real waiting


def test_delivery_dead_letters_after_max_attempts() -> None:
    transport = FakeTransport([500, 500, 500, 500, 500])
    deliverer, sleeps = make_deliverer(transport)
    outcome = deliverer.deliver("https://ops.northwind-robotics.example/hook", make_event())
    assert outcome is DeliveryOutcome.DEAD_LETTERED
    assert sleeps == [1.0, 2.0, 4.0, 8.0]
    assert len(deliverer.dead_letters) == 1
    assert deliverer.dead_letters[0].reason == "max attempts exhausted"


def test_delivery_does_not_retry_permanent_4xx() -> None:
    transport = FakeTransport([422])
    deliverer, sleeps = make_deliverer(transport)
    outcome = deliverer.deliver("https://ops.northwind-robotics.example/hook", make_event())
    assert outcome is DeliveryOutcome.REJECTED_PERMANENT
    assert len(transport.calls) == 1 and sleeps == []
    assert deliverer.dead_letters[0].reason == "permanent 422"


def test_network_error_is_retryable() -> None:
    transport = FakeTransport([OSError, 200])
    deliverer, _ = make_deliverer(transport)
    outcome = deliverer.deliver("https://ops.northwind-robotics.example/hook", make_event())
    assert outcome is DeliveryOutcome.DELIVERED
    assert transport.calls[0][1]["x-northwind-signature"].startswith(f"t={NOW},v1=")
\end{minted}

\subsection{Server-Sent Events: Producer, Parser, and Resume}
\label{sec:ch11-code-sse}

Listing~11.4 is both halves of the SSE contract. The producer is a FastAPI
\texttt{StreamingResponse}~\cite{fastapidocs} emitting well-formed
\texttt{text/event-stream} --- \texttt{event}/\texttt{data}/\texttt{id}/
\texttt{retry} fields, multi-line data folding, comment keep-alives, the
\texttt{X-Accel-Buffering: no} escape hatch --- and honoring
\texttt{Last-Event-ID} resume from the scripted feed. The client parser is a
pure generator over lines implementing the WHATWG dispatch rules, so it works
over \texttt{httpx}~\cite{httpxdocs}, the stdlib, or a test fixture.

\begin{minted}{python}
"""Server-Sent Events: FastAPI producer + a dependency-free client parser."""
from __future__ import annotations

from collections.abc import AsyncIterator, Iterable, Iterator
from dataclasses import dataclass

from fastapi import FastAPI, Request
from fastapi.responses import StreamingResponse

# The scripted Northwind Robotics telemetry feed (synthetic, deterministic).
TELEMETRY_FEED: list[tuple[int, str, str]] = [
    (1, "telemetry", '{"robot": "NWR-07", "battery": 0.82}'),
    (2, "telemetry", '{"robot": "NWR-07", "battery": 0.79}'),
    (3, "alert", '{"robot": "NWR-07", "alert": "dock requested"}'),
]
RETRY_MS = 2000


def format_sse(data: str, *, event: str | None = None,
               event_id: int | None = None, retry: int | None = None) -> str:
    """Serialize one SSE event block. Multi-line data gets one "data:" per line."""
    lines: list[str] = []
    if event is not None:
        lines.append(f"event: {event}")
    for data_line in data.splitlines():
        lines.append(f"data: {data_line}")
    if event_id is not None:
        lines.append(f"id: {event_id}")
    if retry is not None:
        lines.append(f"retry: {retry}")
    return "\n".join(lines) + "\n\n"


@dataclass(frozen=True)
class ServerSentEvent:
    data: str
    event: str = "message"
    event_id: str | None = None
    retry: int | None = None


def parse_sse(lines: Iterable[str]) -> Iterator[ServerSentEvent]:
    """Parse text/event-stream lines per the WHATWG dispatch rules.

    A blank line dispatches the buffered event; lines starting with ":" are
    comments (keep-alives) and are ignored.
    """
    data_parts: list[str] = []
    event_type: str | None = None
    event_id: str | None = None
    retry: int | None = None

    def dispatch() -> ServerSentEvent | None:
        nonlocal data_parts, event_type
        if not data_parts:
            event_type = None
            return None
        parsed = ServerSentEvent("\n".join(data_parts),
                                 event_type or "message", event_id, retry)
        data_parts, event_type = [], None
        return parsed

    for raw in lines:
        line = raw.rstrip("\r\n")
        if line == "":
            dispatched = dispatch()
            if dispatched is not None:
                yield dispatched
        elif line.startswith(":"):
            continue
        else:
            field, _, value = line.partition(":")
            if value.startswith(" "):
                value = value[1:]
            if field == "data":
                data_parts.append(value)
            elif field == "event":
                event_type = value
            elif field == "id":
                event_id = value
            elif field == "retry":
                if value.isdigit():
                    retry = int(value)
    dispatched = dispatch()
    if dispatched is not None:
        yield dispatched


def create_app() -> FastAPI:
    app = FastAPI(title="Northwind Robotics Telemetry Stream")

    async def event_generator(last_event_id: int) -> AsyncIterator[str]:
        # Announce the reconnection delay first, then resume AFTER last_event_id.
        yield format_sse("connected", event="control", retry=RETRY_MS)
        for event_id, event_type, data in TELEMETRY_FEED:
            if event_id > last_event_id:
                yield format_sse(data, event=event_type, event_id=event_id)
        yield ": keep-alive\n\n"

    @app.get("/stream")
    async def stream(request: Request) -> StreamingResponse:
        raw_last_id = request.headers.get("last-event-id", "0")
        last_event_id = int(raw_last_id) if raw_last_id.isdigit() else 0
        return StreamingResponse(
            event_generator(last_event_id),
            media_type="text/event-stream",
            headers={"Cache-Control": "no-cache", "X-Accel-Buffering": "no"},
        )

    return app
\end{minted}

The integration tests (Listing~11.5) exercise the full stack with FastAPI's
in-process \texttt{TestClient} --- no sockets --- and prove the resume
contract: reconnecting with \texttt{Last-Event-ID: 2} replays only event 3.

\begin{minted}{python}
"""Offline tests for the SSE producer, parser, and Last-Event-ID resume."""
from __future__ import annotations

from fastapi.testclient import TestClient

from sse_streams import RETRY_MS, create_app, format_sse, parse_sse

FIXTURE_STREAM = (
    ": this is a comment keep-alive\n"
    "\n"
    "event: telemetry\n"
    "data: {\"robot\": \"NWR-07\",\n"
    "data: \"battery\": 0.82}\n"
    "id: 42\n"
    "retry: 1500\n"
    "\n"
    "data: bare message with default event type\n"
    "\n"
)


def test_parser_handles_comments_multiline_data_id_and_retry() -> None:
    events = list(parse_sse(FIXTURE_STREAM.splitlines(keepends=True)))
    assert len(events) == 2
    first, second = events
    assert first.event == "telemetry"
    assert first.data == '{"robot": "NWR-07",\n"battery": 0.82}'  # multi-line join
    assert first.event_id == "42"
    assert first.retry == 1500
    assert second.event == "message"  # default type
    assert second.retry == 1500  # retry persists until replaced


def test_format_sse_round_trips_through_parser() -> None:
    wire = format_sse("line one\nline two", event="alert", event_id=7, retry=500)
    (event,) = list(parse_sse(wire.splitlines(keepends=True)))
    assert (event.data, event.event, event.event_id, event.retry) == (
        "line one\nline two", "alert", "7", 500,
    )


def test_server_streams_all_events_with_ids() -> None:
    client = TestClient(create_app())
    with client.stream("GET", "/stream") as response:
        assert response.status_code == 200
        assert response.headers["content-type"].startswith("text/event-stream")
        assert response.headers["x-accel-buffering"] == "no"
        events = list(parse_sse(response.iter_lines()))
    ids = [e.event_id for e in events if e.event_id is not None]
    assert ids == ["1", "2", "3"]
    assert events[0].retry == RETRY_MS
    assert events[-1].event == "alert"


def test_server_resumes_after_last_event_id() -> None:
    """Reconnect with Last-Event-ID: 2 -> only event 3 is replayed."""
    client = TestClient(create_app())
    with client.stream("GET", "/stream", headers={"Last-Event-ID": "2"}) as response:
        events = list(parse_sse(response.iter_lines()))
    replayed = [e for e in events if e.event_id is not None]
    assert [e.event_id for e in replayed] == ["3"]
    assert replayed[0].event == "alert"
\end{minted}

\subsection{WebSockets: Handshake Math and an Echo/Broadcast Demo}
\label{sec:ch11-code-ws}

Listing~11.6 proves the handshake and then drives a real (localhost-only)
WebSocket session with the \texttt{websockets} library: echo to the sender,
broadcast to peers, protocol-level ping/pong, and a clean 1000 close.
Synchronization is via \texttt{asyncio.Event}, never \texttt{sleep}-based
timing assertions.

\begin{minted}{python}
"""Minimal WebSocket echo/broadcast demo + the RFC 6455 handshake math."""
from __future__ import annotations

import asyncio
import base64
import hashlib
from typing import Any

from websockets.asyncio.client import connect
from websockets.asyncio.server import ServerConnection, serve

WS_GUID = "258EAFA5-E914-47DA-95CA-C5AB0DC85B11"  # fixed by RFC 6455 section 1.3


def sec_websocket_accept(key: str) -> str:
    """Derive Sec-WebSocket-Accept from the client's Sec-WebSocket-Key."""
    sha1 = hashlib.sha1((key + WS_GUID).encode("ascii")).digest()
    return base64.b64encode(sha1).decode("ascii")


class BroadcastHub:
    """Tracks connected robot consoles; every message is echoed and broadcast."""

    def __init__(self) -> None:
        self.clients: set[ServerConnection] = set()

    async def handler(self, ws: ServerConnection) -> None:
        self.clients.add(ws)
        try:
            async for message in ws:
                await ws.send(f"echo:{message}")  # direct echo to the sender
                for peer in self.clients - {ws}:
                    await peer.send(f"broadcast:{message}")
        finally:
            self.clients.discard(ws)


async def run_server(hub: BroadcastHub) -> Any:
    return await serve(hub.handler, "127.0.0.1", 0)  # port 0 => OS picks a free port


async def demo() -> list[str]:
    """Deterministic scripted session; asyncio.Event provides synchronization."""
    hub = BroadcastHub()
    received: list[str] = []
    peer_ready = asyncio.Event()

    async with await run_server(hub) as server:
        port = server.sockets[0].getsockname()[1]

        async def peer() -> None:
            async with connect(f"ws://127.0.0.1:{port}") as ws:
                peer_ready.set()
                async for message in ws:
                    received.append(f"peer got {message}")

        peer_task = asyncio.create_task(peer())
        await peer_ready.wait()  # explicit synchronization, never sleep-based

        async with connect(f"ws://127.0.0.1:{port}") as ws:
            await ws.send("NWR-07 online")
            received.append(await ws.recv())  # our own echo

            pong = await ws.ping()  # protocol-level keep-alive (opcode 0x9/0xA)
            await asyncio.wait_for(pong, timeout=5)
            received.append("pong received")

            await ws.close(code=1000, reason="demo complete")  # clean close
        peer_task.cancel()

    return received


if __name__ == "__main__":
    print(sec_websocket_accept("dGhlIHNhbXBsZSBub25jZQ=="))  # RFC 6455 example
    for line in asyncio.run(demo()):
        print(line)
\end{minted}

Running \texttt{python ws\_demo.py} prints
\texttt{s3pPLMBiTxaQ9kYGzzhZRbK+xOo=} --- the exact accept value from the
RFC~6455 worked example --- followed by the echo, the peer's broadcast
receipt, and the pong. The tests (Listing~11.7) assert all of it.

\begin{minted}{python}
"""Offline tests for the WebSocket demo and the RFC 6455 handshake math."""
from __future__ import annotations

import asyncio

from ws_demo import BroadcastHub, demo, sec_websocket_accept


def test_sec_websocket_accept_matches_rfc6455_example() -> None:
    # Worked example straight from RFC 6455 section 1.3.
    assert sec_websocket_accept("dGhlIHNhbXBsZSBub25jZQ==") == \
        "s3pPLMBiTxaQ9kYGzzhZRbK+xOo="


def test_echo_broadcast_ping_and_clean_close() -> None:
    received = asyncio.run(demo())
    assert "echo:NWR-07 online" in received
    assert "peer got broadcast:NWR-07 online" in received  # delivered to the peer
    assert "pong received" in received


def test_hub_disconnect_cleans_registry() -> None:
    async def scenario() -> BroadcastHub:
        from websockets.asyncio.client import connect
        from websockets.asyncio.server import serve

        hub = BroadcastHub()
        async with await serve(hub.handler, "127.0.0.1", 0) as server:
            port = server.sockets[0].getsockname()[1]
            async with connect(f"ws://127.0.0.1:{port}"):
                assert len(hub.clients) == 1
            # allow the server task to run its finally block deterministically
            for _ in range(10):
                await asyncio.sleep(0)
                if not hub.clients:
                    break
        return hub

    hub = asyncio.run(scenario())
    assert hub.clients == set()
\end{minted}

\subsection{The Event Contract in AsyncAPI}
\label{sec:ch11-code-asyncapi}

Listing~11.8 is the AsyncAPI~3.0 description of the same Northwind Robotics
telemetry contract~\cite{asyncapi}. Note what the document buys you beyond
prose: the channel \emph{address} is a templated topic string with named
parameters; \emph{operations} declare direction explicitly (\texttt{send}
from the fleet's perspective, \texttt{receive} for consumers); message
payloads carry JSON Schemas with patterns and ranges, so a CI gate can reject
a producer that starts emitting \texttt{battery\_level: 12}; and the
\texttt{servers} block documents the binding (Kafka, port 9092) without the
logical contract depending on it. The same document drives code generation,
documentation portals, and contract-diff gates --- the OpenAPI toolchain's
event-driven mirror.

\begin{minted}{yaml}
asyncapi: "3.0.0"
info:
  title: Northwind Robotics Telemetry Events
  version: "1.0.0"
  description: >-
    Synthetic event contract for the Northwind Robotics fleet. Consumed by
    the operations dashboard (SSE) and by customer webhook endpoints.
servers:
  production:
    host: broker.northwind-robotics.example:9092
    protocol: kafka
    description: Fictional in-cluster broker (documentation only; tests stay offline).
channels:
  robotTelemetry:
    address: northwind.robots.{robot_id}.telemetry
    messages:
      telemetryReading:
        $ref: "#/components/messages/TelemetryReading"
      batteryAlert:
        $ref: "#/components/messages/BatteryAlert"
operations:
  publishTelemetry:
    action: send
    channel:
      $ref: "#/channels/robotTelemetry"
  consumeTelemetry:
    action: receive
    channel:
      $ref: "#/channels/robotTelemetry"
components:
  messages:
    TelemetryReading:
      payload:
        type: object
        required: [id, robot_id, battery_level, recorded_at]
        properties:
          id: { type: string, pattern: "^evt_[0-9]+$" }
          robot_id: { type: string, pattern: "^NWR-[0-9]{2}$" }
          battery_level: { type: number, minimum: 0, maximum: 1 }
          recorded_at: { type: string, format: date-time }
    BatteryAlert:
      payload:
        type: object
        required: [id, robot_id, threshold, recorded_at]
        properties:
          id: { type: string }
          robot_id: { type: string }
          threshold: { type: number, minimum: 0, maximum: 1 }
          recorded_at: { type: string, format: date-time }
\end{minted}

\subsection{Contrast: Long Polling in Twelve Lines}
\label{sec:ch11-code-longpoll}

Listing~11.9 is the long-poll alternative to Listing~11.4, kept deliberately
small. The client issues \texttt{GET /notifications/next?after=N}; the server
holds the request until a notification is published or a server-side deadline
(always below the proxy idle timeout) expires, then answers a normal JSON
document --- possibly \texttt{null} --- and the client immediately re-asks.
Every request is an ordinary, cacheable-in-principle, auth-compatible HTTP
exchange; the cost is a held worker per waiting client and a reconnect after
every single event.

\begin{minted}{python}
"""Long polling, for contrast with SSE: hold the request until data or deadline."""
from __future__ import annotations

import asyncio

from fastapi import FastAPI, Query
from pydantic import BaseModel

HOLD_SECONDS = 15.0  # server-side deadline, always below the proxy idle timeout


class Notification(BaseModel):
    notification_id: int
    robot: str
    message: str


class NotificationBus:
    def __init__(self) -> None:
        self._waiters: list[asyncio.Queue[Notification]] = []

    def subscribe(self) -> asyncio.Queue[Notification]:
        queue: asyncio.Queue[Notification] = asyncio.Queue()
        self._waiters.append(queue)
        return queue

    def publish(self, note: Notification) -> None:
        for queue in self._waiters:
            queue.put_nowait(note)


def create_app(bus: NotificationBus) -> FastAPI:
    app = FastAPI(title="Northwind Robotics Long-Poll Notifications")

    @app.get("/notifications/next", response_model=Notification | None)
    async def next_notification(after: int = Query(0, ge=0)) -> Notification | None:
        queue = bus.subscribe()
        try:
            return await asyncio.wait_for(queue.get(), timeout=HOLD_SECONDS)
        except TimeoutError:
            return None  # 200 with null body: client reconnects immediately

    return app
\end{minted}

\begin{examplebox}
\textbf{Execution evidence.} All listings in this section were executed
offline on Python~3.13 (Windows, PowerShell) before inclusion:
\texttt{pytest -q} reports \textbf{18 passed} --- eleven webhook tests
(tamper, stale timestamp, wrong secret, replay idempotency, rotation-header
parsing, delivery success, backoff schedule, DLQ exhaustion, permanent-4xx
short-circuit, transport-error retry), four SSE tests (parser grammar,
round-trip, full stream, \texttt{Last-Event-ID} resume), and three WebSocket
tests (RFC~6455 accept vector, echo/broadcast/ping/close, registry cleanup).
\texttt{python ws\_demo.py} prints the RFC accept value
\texttt{s3pPLMBiTxaQ9kYGzzhZRbK+xOo=} followed by the scripted session
transcript. The AsyncAPI document parses cleanly as YAML~3.0.
\end{examplebox}

%% ----------------------------------------------------------------------------
\section{Common Pitfalls \& Anti-Patterns}
\label{sec:ch11-pitfalls}

Each entry below is a production scar, not a hypothetical.

\subsection{Unsigned webhooks}
The single most exploited webhook bug. An endpoint with no signature check is
an unauthenticated ``do anything'' API reachable by anyone who scrapes the
URL from client-side code, DNS logs, or a support forum post. Even a
configured-but-unenforced check (``verify if the header is present'') fails:
attackers simply omit the header. The fix is fail-closed verification ---
missing, malformed, stale, or mismatched all reject --- as in Listing~11.1.

\subsection{Retrying permanent failures forever}
A delivery worker that retries 4xx responses indistinguishably from 5xx turns
every consumer bug into an amplification loop: the consumer deploys a schema
change, starts rejecting with 422, and the provider hammers it at full retry
rate for days. Bound the budget, separate retryable from permanent statuses,
and dead-letter the rest (Listing~11.2). The mirror-image consumer bug is
returning 500 for a \emph{payload} problem --- which invites retries that can
never succeed --- instead of a crisp 4xx.

\subsection{Assuming ordering}
Consumers that apply events blindly in arrival order corrupt state under
retries and parallelism: \texttt{robot.updated} (battery 0.79) arrives after
\texttt{robot.updated} (battery 0.82) and the dashboard regresses. Either
send full state per event and make consumers last-write-wins on a
higher-is-newer cursor (\texttt{occurred\_at} or a sequence number), or
buffer and reorder. ``Our provider delivers in order'' is a claim about
today's implementation, not the contract.

\subsection{SSE through buffering proxies}
The stream works on the developer's machine and stalls in production because
nginx's default \texttt{proxy\_buffering on} accumulates events before
forwarding. Symptoms: clients receive nothing for minutes, then a burst.
Defenses: \texttt{X-Accel-Buffering: no}, \texttt{Cache-Control: no-cache},
flush per event, comment keep-alives, and --- the real fix --- an end-to-end
test that runs \emph{through} the production proxy path, not around it.

\subsection{WebSocket for plain CRUD}
Rebuilding \texttt{POST /robots} as a WebSocket message throws away caching,
idempotency keys, status codes, tracing, and every gateway feature your
organization already paid for --- and adds connection lifecycle, reconnection,
and backpressure code you now own. Bidirectional channels earn their keep
only when both directions carry meaningful, latency-sensitive traffic.

\subsection{Non-idempotent receivers}
Processing a webhook before acknowledging it --- \texttt{charge\_card()} then
return 200 --- means every lost response re-charges the customer. Record the
event id transactionally with the effect, or queue-then-ack, so redelivery is
a no-op. Note the ordering subtlety: ack \emph{before} durable recording and
a crash loses the event; process \emph{before} durable recording and a crash
repeats it. Idempotency is what makes the second ordering safe.

\subsection{Leaking event payloads in logs}
Webhook payloads routinely carry PII-adjacent operational data; delivery
services that log full bodies at INFO, or receivers that log the raw request
``for debugging,'' build an ungoverned replica of the event stream inside the
log pipeline. Log event ids, types, sizes, and outcomes --- never bodies,
never signature secrets, never full headers.

\subsection{Sleep-based async tests}
\texttt{await asyncio.sleep(0.3)} as a synchronization mechanism makes tests
flaky under CI load and slow when generous. Use \texttt{asyncio.Event},
queues, and \texttt{wait\_for} with explicit timeouts
(Listings~11.6--11.7): the test passes the instant the condition holds and
fails loudly when it never will.

%% ----------------------------------------------------------------------------
\section{Hands-On Production Lab \& Real-World Case Study}
\label{sec:ch11-lab}

\subsection{Lab: The Northwind Robotics Real-Time Notifications Service}

You will assemble this chapter's components into one coherent service: a
\emph{notification fan-out} for the Northwind Robotics fleet that (a) serves
an SSE feed to browser dashboards and (b) delivers signed webhooks to two
registered customer endpoints, one of which is deliberately flaky.

\textbf{Setup} (PowerShell):

\begin{minted}{text}
python -m venv .venv
.\.venv\Scripts\Activate.ps1
pip install fastapi httpx uvicorn pydantic websockets pytest
mkdir lab_ch11; cd lab_ch11
\end{minted}

\textbf{Step 1 --- the receiver.} Save Listing~11.1 as
\texttt{webhook\_receiver.py}. This is the \emph{customer's} endpoint;
you will run two copies in-process via \texttt{TestClient}: one configured
with the correct secret, one misconfigured (wrong secret) to act as the
flaky endpoint.

\textbf{Step 2 --- the delivery worker.} Save Listing~11.2 as
\texttt{webhook\_delivery.py}. Replace \texttt{FakeTransport} with a real
transport built on \texttt{httpx}~\cite{httpxdocs}: a function
\texttt{(url, headers, body) -> int} that POSTs with a 5\,s timeout and
returns the status code, translating \texttt{httpx.TransportError} to
\texttt{OSError}. Register it against the healthy receiver's URL and observe
\texttt{DELIVERED}; register the misconfigured one and observe
\texttt{REJECTED\_PERMANENT} (the receiver's 400) with exactly one attempt.

\textbf{Step 3 --- the SSE feed.} Save Listing~11.4 as
\texttt{sse\_streams.py} and serve it with the app factory:
\texttt{uvicorn "sse\_streams:create\_app" {-}{-}factory}. In a second
terminal, consume it with \texttt{curl -N http://127.0.0.1:8000/stream} and
watch events arrive with ids; kill and re-run curl with
\texttt{-H "Last-Event-ID: 2"} and confirm only event 3 replays.

\textbf{Step 4 --- wire the fan-out.} Write a thin orchestrator: each event
in \texttt{TELEMETRY\_FEED} is both appended to the SSE feed \emph{and}
handed to the \texttt{WebhookDeliverer} for each registered endpoint. Run the
deterministic test suite (Listings~11.3, 11.5) as the acceptance gate.

\textbf{Acceptance checks:} healthy endpoint receives all three events with
valid signatures; misconfigured endpoint is dead-lettered with reason
\texttt{permanent 400}; SSE clients resume correctly; \texttt{pytest -q} is
green with no network access (unplug it and prove it).

\subsection{War Story 1: The Retry Loop That DDoSed a Customer}

A fleet-management provider (the fictionalized composite is Northwind
Robotics, the pattern is industry-common) shipped webhook delivery with a
single flaw: the retry predicate was \texttt{status != 200}. A customer's
receiver began rejecting deliveries with 422 after a schema change on their
side. The provider's worker retried --- immediately, then on a short fixed
interval, with no cap. Queued events accumulated during the customer's
maintenance window; when the window closed, forty thousand queued deliveries
retried in near-synchrony against an endpoint that rejected each one in
milliseconds. From the customer's load balancer the traffic was
indistinguishable from an application-layer DDoS; their WAF rate-limited the
provider's egress IPs, which made \emph{every} customer's deliveries to that
domain fail, which queued more events, which deepened the storm.

The postmortem produced exactly the rules encoded in Listing~11.2: classify
statuses (permanent 4xx $\Rightarrow$ stop immediately), bound attempts,
exponential backoff \emph{with jitter} so recovering queues do not fire in
lockstep, per-endpoint circuit breaking, and a DLQ with paging alerts. The
metric that would have caught it in minutes --- deliveries-per-endpoint-per-
minute and 4xx ratio --- became a standing dashboard.

\subsection{War Story 2: The Unsigned-Webhook Spoofing Incident}

A robotics-operations startup exposed \texttt{POST /hooks/fleet} to receive
``maintenance completed'' events from its CMMS vendor. The endpoint was
``protected'' by obscurity: a long random path. A former contractor, who knew
the path from integration work, POSTed a fabricated \texttt{maintenance.
completed} event for a robot that was, in reality, still faulted; the
dispatcher cleared the work order automatically and the robot was scheduled
back into a production shift. The resulting equipment damage was modest; the
audit finding was not --- the endpoint had accepted a forged state transition
with zero authentication, and logs showed the vendor's ``integration guide''
had suggested signature verification as \emph{optional}.

The remediation was Listing~11.1 verbatim: shared-secret HMAC over
timestamp-plus-body, a 300-second tolerance window, constant-time comparison,
idempotency on event id, and --- organizationally --- a policy that no
webhook endpoint merges without a red-team test that submits an unsigned
payload and asserts rejection. Obscurity of a URL is not a control; it is a
hope.

%% ----------------------------------------------------------------------------
\section{Self-Assessment Exercises \& Deep-Dive Questions}
\label{sec:ch11-exercises}

\begin{enumerate}[label=\textbf{E11.\arabic*.}, leftmargin=2.2cm]
  \item \textbf{Pattern triage.} For each requirement, pick exactly one
        pattern from Table~\ref{tab:ch11-pattern-matrix} and defend the pick:
        (a) a stock-ticker widget on a public marketing page; (b) a
        two-player browser chess game; (c) ``notify our ERP when an invoice
        is approved''; (d) a mobile app's badge count on a metered
        connection; (e) a SOC wallboard that must survive the corporate
        proxy.
  \item \textbf{Signature forensics.} A receiver logs show a burst of
        deliveries with valid signatures but timestamps 40 minutes old. All
        were rejected. What two explanations fit, and what one-line config
        change on the provider side would distinguish them?
  \item \textbf{Rotation drill.} Extend Listing~11.2 so the deliverer signs
        each attempt with \emph{two} secrets during a rotation window
        (\texttt{t=...,v1=...,v1=...}). Prove with a test that receivers
        holding either secret accept.
  \item \textbf{SSE grammar fuzz.} The parser in Listing~11.4 ignores fields
        it does not know. Per the WHATWG rules, what should happen for:
        a line with no colon; a \texttt{data} field with an empty value; an
        \texttt{id} containing a NUL; \texttt{retry: abc}? Encode your
        answers as four new tests.
  \item \textbf{Frame dissection.} By hand, write the two header bytes of:
        an unmasked final text frame with a 60-byte payload; a masked final
        binary frame with a 300-byte payload; an unfragmented ping with
        payload \texttt{"hi"}. Then verify byte two's mask bit in each.
  \item \textbf{Masking math.} Client payload byte \texttt{0x48} (`H'),
        masking-key byte \texttt{0x37}. What goes on the wire? Why could the
        attacker-in-the-browser not simply choose payload bytes that survive
        masking unchanged?
  \item \textbf{Backoff analysis.} Listing~11.2 uses delays
        $\min(b \cdot 2^{n-1}, 60)$. For $b=1$, $n_{\max}=8$: total worst-
        case wait? Now add $\pm 20\%$ jitter: what failure mode does jitter
        prevent, and why is \texttt{random} seeded in the test suite if you
        implement it?
  \item \textbf{Resume-window design.} Your SSE server keeps a 500-event ring
        buffer. A dashboard disconnects for an hour during which 2\,000
        events pass. Design the reconnection behavior: what does the server
        send, what should the client do, and how does the user experience
        degrade gracefully?
  \item \textbf{DLQ triage runbook.} Write the on-call runbook entry for the
        alert ``dead-letter rate $>$ 10/hour'': first three diagnostic
        queries, the decision tree for replay-versus-discard, and the
        customer-communication trigger point.
  \item \textbf{Contract gate.} Add a CI check (a short Python script) that
        loads Listing~11.8 and fails the build if any message payload
        property is removed or retyped relative to a checked-in previous
        version. Which compatibility mode does your check implement?
\end{enumerate}

%% ----------------------------------------------------------------------------
\section{Interview Q\&A Vault}
\label{sec:ch11-interview}

\paragraph{Junior tier}

\begin{enumerate}[label=\textbf{Q\arabic*.}, leftmargin=1.6cm]
  \item \textbf{What is the difference between polling and webhooks?}
        Polling: the consumer repeatedly asks ``anything new?'' on its own
        schedule --- simple, but most responses are empty and latency is
        bounded by the interval. A webhook inverts it: the consumer registers
        a URL and the \emph{provider} POSTs each event as it happens.
        Polling costs requests; webhooks cost you a public, secured endpoint.
  \item \textbf{What does SSE stand for and what is it?}
        Server-Sent Events: a server responds to a GET with
        \texttt{Content-Type: text/event-stream} and keeps the response open,
        appending \texttt{field: value} event blocks separated by blank
        lines. One-directional, server-to-client, text-only, native in
        browsers via \texttt{EventSource}.
  \item \textbf{Name the four SSE fields and what each does.}
        \texttt{data} (payload; multiple lines join with newlines),
        \texttt{event} (named type; default \texttt{message}),
        \texttt{id} (cursor used for \texttt{Last-Event-ID} resume),
        \texttt{retry} (client reconnection delay in ms). Colon-prefixed
        lines are comment keep-alives.
  \item \textbf{What HTTP status begins a WebSocket connection, and what
        headers matter?}
        \texttt{101 Switching Protocols}. The client sends
        \texttt{Upgrade: websocket}, \texttt{Connection: Upgrade},
        \texttt{Sec-WebSocket-Key}, \texttt{Sec-WebSocket-Version: 13}; the
        server replies with \texttt{Sec-WebSocket-Accept}.
  \item \textbf{Is SSE one-way or two-way? What about WebSocket?}
        SSE is strictly server-to-client; clients respond over ordinary HTTP
        requests if needed. WebSocket is full-duplex: both peers send frames
        at any time after the upgrade.
  \item \textbf{Why do webhooks usually sign their payloads?}
        The receiver's URL is public; anyone who finds it could POST forged
        events. An HMAC over the body with a shared secret proves the sender
        holds the secret and the body is unaltered.
  \item \textbf{What is an idempotent webhook receiver and why does it need
        to be one?}
        One where processing the same event twice has the same effect as
        once --- typically by recording the event id and skipping duplicates.
        Needed because delivery is at-least-once: providers retry whenever
        the 2xx is lost, so duplicates are guaranteed.
  \item \textbf{What happens when an SSE connection drops?}
        \texttt{EventSource} waits the last \texttt{retry} interval (default
        3\,s), reconnects, and sends \texttt{Last-Event-ID} with the last
        received \texttt{id} so the server can resume the stream.
  \item \textbf{Give two close codes you would actually use and their
        meanings.}
        1000 normal closure (work complete); 1001 going away (server restart
        or page navigation). Mentioning that 1006 is reserved for
        \emph{abnormal} closure and must never be sent earns extra credit.
  \item \textbf{Webhooks versus a message queue --- when which?}
        Webhooks cross organizational boundaries over public HTTP with zero
        consumer infrastructure beyond an endpoint. Queues (Kafka, SQS) are
        inside one trust/domain boundary, give stronger ordering/retention
        tooling, and require consumer connectivity to the broker.
\end{enumerate}

\paragraph{Mid tier}

\begin{enumerate}[label=\textbf{Q\arabic*.}, leftmargin=1.6cm, start=11]
  \item \textbf{Derive \texttt{Sec-WebSocket-Accept}.}
        Concatenate the client's \texttt{Sec-WebSocket-Key} with the fixed
        GUID \texttt{258EAFA5-...-C5AB0DC85B11}, take SHA-1, Base64-encode.
        RFC vector: \texttt{dGhlIHNhbXBsZSBub25jZQ==} $\rightarrow$
        \texttt{s3pPLMBiTxaQ9kYGzzhZRbK+xOo=}. Purpose: prove the server
        implemented WebSocket, not just header echo.
  \item \textbf{Why are client-to-server WebSocket frames masked?}
        Not confidentiality --- intermediaries. A hostile browser script
        could otherwise shape outbound bytes into something a misbehaving
        transparent proxy parses as a smuggled HTTP request, poisoning shared
        caches. A per-frame random 32-bit XOR key chosen \emph{after} the
        app supplies bytes makes the wire bytes unpredictable to the
        attacker.
  \item \textbf{Walk me through the Stripe-style \texttt{t=...,v1=...}
        verification.}
        Parse timestamp and signature(s); check $|now - t| \le$ tolerance
        (replay window); recompute HMAC-SHA256 over
        ``$t$~\texttt{.}~raw-body'' with the endpoint secret; constant-time
        compare against every \texttt{v1}; only then parse the JSON. Verify
        raw bytes because re-serialization is not byte-stable.
  \item \textbf{Why multiple \texttt{v1=} values in one header?}
        Secret rotation: during the overlap window the provider signs with
        old and new secrets; receivers configured with either still verify,
        so rotation needs no synchronized flag day.
  \item \textbf{Design webhook delivery guarantees.}
        Durable queue per endpoint; at-least-once delivery with unique event
        ids; retry only retryable statuses (5xx, 429, timeouts, connection
        errors) with capped exponential backoff plus jitter; permanent 4xx
        dead-letters immediately; DLQ with alerting and replay; document
        best-effort ordering; consumer side: idempotent processing on event
        id.
  \item \textbf{SSE versus WebSockets for a live operations dashboard ---
        decide.}
        Dashboards are one-way server-to-client: SSE. You get automatic
        reconnection with \texttt{Last-Event-ID} resume, plain-HTTP infra
        compatibility (CORS, cookies, auth), trivial debugging with
        \texttt{curl}. Choose WebSocket only if operators must also
        \emph{send} low-latency commands on the same channel.
  \item \textbf{Your SSE stream works locally but not behind nginx. Diagnose.}
        Proxy buffering. Fix with \texttt{X-Accel-Buffering: no},
        \texttt{Cache-Control: no-cache}, per-event flush, comment
        keep-alives every 10--15\,s. Verify through the real proxy path.
  \item \textbf{HTTP/1.1 versus HTTP/2 for SSE?}
        HTTP/1.1: browsers cap at six connections per origin; two
        long-lived streams can starve other requests --- use a separate
        origin or HTTP/2. HTTP/2 multiplexes many streams over one
        connection, removing the cap, at the cost of shared TCP head-of-line
        blocking.
  \item \textbf{How does \texttt{Last-Event-ID} resume work, and what is its
        limit?}
        The client replays the last seen \texttt{id} on reconnect; the
        server replays events after it from a buffer/log. The limit is the
        retention window: if the client fell further behind than the buffer,
        the server must signal a resync (control event or 204) and the client
        re-fetches current state.
  \item \textbf{Exactly-once webhook delivery --- possible?}
        No; over an unreliable network the sender cannot distinguish
        ``receiver crashed before processing'' from ``processed but ack
        lost'' (Two Generals). Honest design: at-least-once transport plus
        idempotent consumers for effectively-once processing.
  \item \textbf{Why does SSE get CORS/cookie auth ``for free'' while
        \texttt{EventSource} cannot set custom headers?}
        SSE \emph{is} an HTTP response, so CORS and cookies apply natively;
        but the \texttt{EventSource} API exposes no header API, so bearer
        tokens must ride as cookies or (carefully) query
        parameters~\cite{rfc6749}. WebSocket's constructor accepts
        subprotocols but not arbitrary headers either; token-in-first-message
        or ticket endpoints are the usual patterns.
  \item \textbf{What is a subprotocol and how is it chosen?}
        An application-level protocol layered on WebSocket framing (e.g.\
        \texttt{mqtt}, \texttt{graphql-ws}), offered by the client in
        \texttt{Sec-WebSocket-Protocol} as a preference list; the server
        picks at most one and returns it; a client that required one must
        close if the pick is absent.
\end{enumerate}

\paragraph{Senior tier}

\begin{enumerate}[label=\textbf{Q\arabic*.}, leftmargin=1.6cm, start=23]
  \item \textbf{Dissect a WebSocket frame header.}
        Byte 0: FIN bit, RSV1--3, 4-bit opcode. Byte 1: MASK bit, 7-bit
        length. Length 0--125 literal; 126 $\Rightarrow$ next 2 bytes are a
        uint16; 127 $\Rightarrow$ next 8 bytes uint63 (big-endian). Then a
        4-byte masking key iff MASK. Control frames (8/9/A) must be FIN with
        $\le$125-byte payloads and may interleave between data fragments.
  \item \textbf{Design a webhook system for one million endpoints.}
        Partitioned durable queue (per-endpoint FIFO only if ordering is
        promised); worker pool with per-endpoint concurrency limits and
        circuit breakers; status classification; capped jittered backoff;
        DLQ per endpoint with replay tooling; signed delivery with rotation;
        endpoint verification at registration; SSRF defenses (egress proxy,
        DNS-pinning, private-range blocking); metrics: attempts, latency,
        4xx/5xx ratio, queue depth per endpoint.
  \item \textbf{How do you scale SSE/WebSocket fan-out horizontally?}
        Connections pin to a process, so introduce a backplane: each node
        subscribes to a broker (Kafka/Redis pub/sub) and forwards to its
        local sockets; sticky routing is unnecessary if every node can serve
        every client; resume state (SSE) must live in the shared log, not
        node memory.
  \item \textbf{Backpressure: a slow SSE consumer and a slow WebSocket
        consumer --- what happens in each, and what do you do?}
        TCP flow control eventually stalls both; the server's send buffer
        grows. SSE: bound a per-client queue, drop with a ``resync'' control
        event (the client can resume via \texttt{Last-Event-ID}). WebSocket:
        monitor buffered amount, apply app-level policies (drop deltas, send
        snapshots, close with 1008/1013); never let unbounded buffering
        become a memory-exhaustion vector.
  \item \textbf{Ordering and deduplication strategy for a webhook consumer
        processing financial events.}
        At-least-once plus possible reordering demands: store event id and
        \texttt{occurred\_at}/sequence transactionally with the effect;
        dedupe on id; make effects last-write-wins on a monotonic cursor or
        buffer-and-reorder within a window; reconcile periodically against
        the provider's pull API (events you \emph{missed} are the silent
        failure mode).
  \item \textbf{Threat-model a public webhook receiver.}
        Forgery (mitigate: HMAC verification), replay (timestamp tolerance +
        idempotency store), tampering (HMAC over raw body), secret leak
        (instant rotation, multiple \texttt{v1}), log injection / PII leak
        (no payload logging), DoS via body size (limit bytes read before
        verifying), timing oracle on compare (constant-time comparison).
  \item \textbf{When would you pick long polling over SSE in 2026?}
        Behind aggressive intermediaries that buffer or kill long-lived
        streams; when every response must be a complete, cacheable HTTP
        document; for very low event rates where held sockets cost more than
        periodic requests; or where the client platform lacks
        \texttt{EventSource} and you refuse a dependency.
  \item \textbf{How do CloudEvents and AsyncAPI complement each other?}
        CloudEvents standardizes the \emph{envelope} of a single event
        (id, source, type, time, subject) so tooling is portable; AsyncAPI
        standardizes the \emph{contract} of the whole API (channels,
        operations, message schemas, bindings). You document channels and
        payloads with AsyncAPI and shape each message's metadata as a
        CloudEvent~\cite{asyncapi}.
  \item \textbf{What is a poison message and how does your pipeline survive
        one?}
        An event that deterministically crashes or rejects every processing
        attempt. Unhandled, it blocks or spins the consumer forever
        (head-of-line). Handling: bounded attempts with backoff, then
        quarantine to a dead-letter topic \emph{with the error context},
        alert, and continue with the next event; replay after a fix.
  \item \textbf{The provider's clock and yours disagree by ten minutes ---
        what breaks and how do you design for it?}
        The freshness check rejects every delivery. Design: tolerance windows
        sized to expected skew (300\,s), NTP discipline on both fleets,
        monitoring on rejection-by-staleness rate, and --- because the
        signature \emph{covers} the timestamp --- attackers cannot simply
        rewrite $t$, so the window is the only knob; keep it tight.
\end{enumerate}

\paragraph{Staff tier}

\begin{enumerate}[label=\textbf{Q\arabic*.}, leftmargin=1.6cm, start=33]
  \item \textbf{Your org runs REST, webhooks, SSE, WebSockets, and gRPC
        streams. What governance keeps this coherent?}
        One contract repository (OpenAPI + AsyncAPI) with CI diff gates;
        one envelope convention (CloudEvents-style id/source/type/time);
        one signing/auth scheme reused across webhook providers; pattern-
        selection review at design time (the taxonomy matrix) so teams
        justify push channels; uniform consumer obligations (idempotency,
        DLQ handling) in a shared library; deprecation policy identical for
        endpoints and event types~\cite{newman2021microservices,asyncapi}.
  \item \textbf{Multi-region webhook delivery: what changes?}
        Deliver from the region nearest the endpoint to cut latency, but the
        queue and DLQ must survive region loss (replicated log); per-endpoint
        state (circuit breakers, backoff cursors) needs global consistency or
        bounded divergence; secrets replicate with the same rotation overlap;
        fail-open regions must not double-deliver \emph{without} the
        consumer's idempotency absorbing it --- which is why the contract
        promises at-least-once even in steady state.
\end{enumerate}

%% ----------------------------------------------------------------------------
\section{External Bibliography \& Further Reading}
\label{sec:ch11-further}

\begin{itemize}
  \item \textbf{AsyncAPI Initiative}, \emph{AsyncAPI Specification 3.0} ---
        the contract format for channels, messages, operations, and bindings
        used in Listing~11.8~\cite{asyncapi}.
  \item \textbf{IETF RFC~6455}, Fette \& Melnikov, \emph{The WebSocket
        Protocol} --- the handshake, the SHA-1/GUID accept derivation, frame
        format, masking rationale, opcodes, and close-code registry.
  \item \textbf{WHATWG}, \emph{Server-Sent Events} (EventSource) living
        standard --- the authoritative \texttt{text/event-stream} parsing
        and reconnection algorithm implemented in Listing~11.4.
  \item \textbf{MDN Web Docs}, \emph{Server-Sent Events / EventSource} ---
        the browser API surface: named event listeners, \texttt{readyState},
        \texttt{lastEventId}, and error handling.
  \item \textbf{Stripe}, \emph{Webhooks documentation} --- the reference
        implementation of the \texttt{t=...,v1=...} HMAC scheme, tolerance
        windows, and endpoint secret rotation this chapter mirrors.
  \item \textbf{IETF RFC~9110}, \emph{HTTP Semantics} --- methods, status
        codes, content negotiation, and the Upgrade mechanism both SSE and
        WebSocket rely on~\cite{rfc9110}.
  \item \textbf{Adam Bellemare}, \emph{Designing Event-Driven Systems}
        (O'Reilly, 2018) --- event log fundamentals, consumer semantics, and
        the data-on-the-outside discipline behind
        Section~\ref{sec:ch11-asyncapi}.
  \item \textbf{CloudEvents (CNCF)}, \emph{CloudEvents Specification 1.0} ---
        the vendor-neutral event envelope (id, source, type, time) worth
        adopting before your event catalog outgrows ad-hoc JSON.
  \item \textbf{Martin Kleppmann}, \emph{Designing Data-Intensive
        Applications} (O'Reilly, 2017) --- the messaging-systems chapters
        behind at-least-once semantics and the Two Generals framing
        \cite{kleppmann2017ddia}.
  \item \textbf{Sam Newman}, \emph{Building Microservices}, 2nd ed.
        (O'Reilly, 2021) --- event-driven collaboration patterns and their
        operational costs~\cite{newman2021microservices}.
  \item \textbf{FastAPI} and \textbf{httpx} documentation ---
        \texttt{StreamingResponse}, \texttt{WebSocket} routes, and streaming
        client usage used throughout Section~\ref{sec:ch11-code}
        \cite{fastapidocs,httpxdocs}.
\end{itemize}

%% ----------------------------------------------------------------------------
\section{Capstone Projects}
\label{sec:ch11-capstones}

\subsection{Capstone 11.1 [junior] --- Signed Webhook Receiver for Robot Alerts}

\textbf{Objective:} Build the consumer half of a webhook contract from
scratch and prove it hostile-input-safe.

\textbf{Inputs/Datasets:} A fixed synthetic feed of 12 Northwind Robotics
alert events (JSON files, provided); one valid signing secret; a tampered
variant of three events; two events signed with a stale timestamp.

\textbf{Steps:}
\begin{enumerate}
  \item Implement the \texttt{t=...,v1=...} verifier (reuse
        Listing~11.1's primitives) with a 300-second tolerance and
        constant-time comparison.
  \item Add an idempotency store keyed on event id with a durability story
        you can explain (a SQLite table is fine).
  \item Write tests for: valid, tampered, stale, wrong-secret, replayed, and
        missing-header deliveries.
\end{enumerate}

\textbf{Acceptance Criteria:}
\begin{itemize}[label=$\square$]
  \item All 12 valid events processed exactly once; all hostile variants
        rejected with 4xx problem-details bodies.
  \item Duplicate deliveries return 200 with \texttt{"status": "duplicate"}
        and cause no side effects.
  \item Test suite passes offline with a fixed injected clock.
\end{itemize}

\textbf{Stretch Goals:} support a two-secret rotation window; emit a
Prometheus-style counter of accepted/rejected/duplicate deliveries.

\subsection{Capstone 11.2 [mid] --- Webhook Delivery Service with DLQ and Replay}

\textbf{Objective:} Build the provider half: queue, bounded jittered retries,
status classification, dead-letter store, and an operator replay command.

\textbf{Inputs/Datasets:} A registry of five synthetic customer endpoints
(simulated by scripted fake transports: healthy, flaky-500, permanent-422,
timeout-then-healthy, always-down); the Capstone~11.1 event feed.

\textbf{Steps:}
\begin{enumerate}
  \item Extend Listing~11.2 with seeded jitter and a persistent DLQ
        (SQLite or JSONL) capturing event, endpoint, and attempt history.
  \item Add a per-endpoint circuit breaker: after five consecutive exhausted
        deliveries, pause the endpoint for a cooldown.
  \item Implement \texttt{replay <event\_id>} re-enqueueing a dead-lettered
        event; prove a replayed event is delivered (and the receiver's
        idempotency absorbs the earlier partial processing).
\end{enumerate}

\textbf{Acceptance Criteria:}
\begin{itemize}[label=$\square$]
  \item The permanent-422 endpoint receives exactly one attempt per event.
  \item The always-down endpoint dead-letters every event; \texttt{replay}
        restores delivery once its transport recovers.
  \item Backoff schedule is asserted deterministically (injected clock and
        sleeper); no test sleeps.
  \item A one-page delivery-semantics document (at-least-once, ordering
        caveats) ships with the service.
\end{itemize}

\textbf{Stretch Goals:} SSRF guardrails on endpoint registration (scheme
allowlist, private-range rejection); per-endpoint delivery metrics endpoint.

\subsection{Capstone 11.3 [mid] --- Resumable SSE Fleet Dashboard Backend}

\textbf{Objective:} Serve a multi-channel SSE feed with correct
\texttt{Last-Event-ID} resume from a bounded retention buffer, and survive a
buffering proxy.

\textbf{Inputs/Datasets:} A deterministic event generator (seeded) emitting
per-robot telemetry for five robots at scripted intervals.

\textbf{Steps:}
\begin{enumerate}
  \item Generalize Listing~11.4: per-robot channels
        (\texttt{/stream?robot=NWR-07}) plus a combined channel, each with
        \texttt{id}, \texttt{retry}, and comment keep-alives every 15
        seconds.
  \item Implement a 500-event ring buffer per channel; resume serves the
        buffer after \texttt{Last-Event-ID}, and answers a
        \texttt{control: resync-required} event when the cursor has scrolled
        out.
  \item Put nginx (or a small buffering reverse proxy you write yourself) in
        front and demonstrate the fix: \texttt{X-Accel-Buffering: no} plus
        flushes, verified by a client that asserts inter-event arrival gaps.
\end{enumerate}

\textbf{Acceptance Criteria:}
\begin{itemize}[label=$\square$]
  \item A client disconnecting and resuming within the retention window
        misses zero events; one resuming outside it receives the resync
        signal.
  \item Through the proxy, events arrive within one second of emission.
  \item Grammar tests cover comments, multi-line data, unknown fields, and
        empty-data blocks.
\end{itemize}

\textbf{Stretch Goals:} HTTP/2 serving; per-channel \texttt{retry} tuning;
EventSource-compatible JavaScript demo page served by the same app.

\subsection{Capstone 11.4 [senior] --- Bidirectional Robot Control Channel over WebSockets}

\textbf{Objective:} Build a production-shaped WebSocket service: subprotocol
negotiation, per-connection auth via first-message ticket, heartbeat policy,
and graceful fleet-wide shutdown.

\textbf{Inputs/Datasets:} The synthetic fleet (five robots); a scripted
command set (dock, undock, set-speed) with per-command acknowledgements;
a chaos harness that kills connections at scripted moments.

\textbf{Steps:}
\begin{enumerate}
  \item Define a subprotocol \texttt{northwind.control.v1} (JSON
        request/response frames with correlation ids) and negotiate it in
        the handshake; reject connections that require an unknown
        subprotocol.
  \item Implement ticket auth: client sends a one-time ticket in the first
        frame; server closes with 1008 on failure.
  \item Server pings every 20\,s; two missed pongs $\Rightarrow$ close 1011
        and reap. Client reconnects with exponential backoff and replays
        unacknowledged commands (deduped server-side by correlation id).
  \item Fleet shutdown: server sends close 1001 (\texttt{going away}) to all
        connections, waits for echoes, and exits before the deploy deadline.
\end{enumerate}

\textbf{Acceptance Criteria:}
\begin{itemize}[label=$\square$]
  \item Commands survive connection kills exactly-once (asserted via the
        chaos harness with seeded scheduling).
  \item Zombie connections are reaped within two heartbeat intervals.
  \item Shutdown completes with zero unclean closes (no 1006) in the happy
        path.
  \item A short design note explains why this use case justifies WebSockets
        over SSE + REST commands.
\end{itemize}

\textbf{Stretch Goals:} permessage-deflate negotiation; a Redis-style
backplane interface (stubbed in-memory) proving the design scales to $N$
server nodes.

\subsection{Capstone 11.5 [staff] --- Event-Contract Governance Platform}

\textbf{Objective:} Stand up the contract-governance layer for Northwind's
entire event surface: AsyncAPI documents, compatibility gates, and a
consumer-obligation library shared by all receivers.

\textbf{Inputs/Datasets:} Three AsyncAPI documents (telemetry, alerts,
maintenance) at two versions each, including one deliberately breaking change
(property retype) and one compatible change (optional property added).

\textbf{Steps:}
\begin{enumerate}
  \item Write an AsyncAPI diff tool classifying channel/message/schema
        changes as breaking or non-breaking (removal, retype, narrowing
        $\Rightarrow$ breaking; additive-optional $\Rightarrow$ safe), exit
        code non-zero on breakage.
  \item Build the shared receiver library: signature verification,
        idempotency store, poison-message quarantine, and structured
        (payload-free) logging --- packaged so Capstone~11.1's receiver can
        adopt it in five lines.
  \item Produce a governance one-pager: compatibility modes, deprecation
        workflow, version-support policy, and the consumer obligations every
        team signs up for.
\end{enumerate}

\textbf{Acceptance Criteria:}
\begin{itemize}[label=$\square$]
  \item The diff gate passes the compatible evolution and fails the breaking
        one, with human-readable diagnostics naming the offending property.
  \item The shared library passes this chapter's full webhook test suite
        unmodified.
  \item The governance document covers schema evolution, secret rotation,
        DLQ ownership, and event-retirement criteria.
\end{itemize}

\textbf{Stretch Goals:} CloudEvents envelope adoption across all three
contracts; a generated Markdown catalog of channels and messages from the
AsyncAPI sources.

%% ----------------------------------------------------------------------------
\section{Python Dependencies Appendix}
\label{sec:ch11-deps}

The chapter's \texttt{requirements.txt}:

\begin{minted}{text}
fastapi>=0.110
websockets>=12
httpx>=0.27
pydantic>=2.7
pytest>=8
\end{minted}

Install under PowerShell (offline machines: pre-stage a wheelhouse with
\texttt{pip download -d wheelhouse -r requirements.txt} on a connected host,
then \texttt{pip install {-}{-}no-index {-}{-}find-links wheelhouse -r
requirements.txt}):

\begin{minted}{text}
python -m venv .venv
.\.venv\Scripts\Activate.ps1
pip install -r requirements.txt
\end{minted}

\subsection{fastapi}
\textbf{Description:} The ASGI framework used for the webhook receiver
(Listing~11.1), the SSE producer (Listing~11.4), and the long-poll contrast
(Listing~11.9), including \texttt{StreamingResponse} and the in-process
\texttt{TestClient} that keeps every test offline~\cite{fastapidocs}.
\textbf{Why included:} it is the book's established server framework
(Chapter~6) and its response/streaming primitives map one-to-one onto the
protocols in this chapter without magic.
\textbf{Alternatives:} Starlette directly (FastAPI's base; fewer batteries);
SSE-Starlette (adds a dedicated \texttt{EventSourceResponse} --- convenient,
but Listing~11.4 deliberately shows the wire format by hand, which is the
pedagogical point); aiohttp or Django Channels if your organization already
standardized there.
\textbf{Installation:} \texttt{pip install "fastapi>=0.110"} in the activated
venv; serving listings interactively additionally needs \texttt{uvicorn}.

\subsection{websockets}
\textbf{Description:} The reference-quality Python WebSocket library;
Listing~11.6 uses its \texttt{websockets.asyncio} implementation (available
since version~12) for both server and client sides of the echo/broadcast
demo, ping/pong, and clean close.
\textbf{Why included:} RFC~6455 correctness (handshake, masking, control
frames, close handshake) is handled for you, letting the chapter focus on
semantics while still exposing the wire details where they matter.
\textbf{Alternatives:} python-socketio (higher-level rooms/namespaces with
automatic fallbacks --- pick it for product chat features, not for learning
the protocol); FastAPI/Starlette's built-in WebSocket routes (fine for
ASGI-hosted endpoints); aiohttp's WS support.
\textbf{Installation:} \texttt{pip install "websockets>=12"}.

\subsection{httpx}
\textbf{Description:} The HTTP client used as the delivery transport in the
lab and as the streaming client paired with the SSE parser; powers
\texttt{TestClient} under the hood~\cite{httpxdocs}.
\textbf{Why included:} one API for sync and async, first-class streaming
response iteration (\texttt{client.stream(...).iter\_lines()}), explicit
timeout configuration --- the three properties webhook delivery and SSE
consumption actually need.
\textbf{Alternatives:} \texttt{urllib.request} (stdlib; workable for the
delivery worker but no ergonomic streaming); \texttt{requests} (no async, no
HTTP/2); \texttt{aiohttp} (async-only).
\textbf{Installation:} \texttt{pip install "httpx>=0.27"}.

\subsection{pydantic}
\textbf{Description:} The validation layer behind FastAPI models (used
directly for the long-poll \texttt{Notification} model in Listing~11.9).
\textbf{Why included:} event payloads are contracts; validating them at the
boundary with typed models is the same discipline Chapter~4 established for
REST representations.
\textbf{Alternatives:} \texttt{dataclasses} plus hand-rolled validation
(what Listings~11.1--11.2 use deliberately, to keep signature logic
dependency-visible); msgspec or attrs for performance-sensitive paths.
\textbf{Installation:} \texttt{pip install "pydantic>=2.7"} (pulled in
transitively by FastAPI).

\subsection{pytest}
\textbf{Description:} The test runner for all eighteen tests in this
chapter.
\textbf{Why included:} fixtures and parametrization keep the hostile-input
matrix (tampered, stale, replayed) declarative; async behavior is tested
through \texttt{asyncio.run} inside synchronous tests, so no pytest-asyncio
plugin is required and scheduling stays explicit.
\textbf{Alternatives:} \texttt{unittest} (stdlib, zero install; more
boilerplate); nose2 (largely superseded).
\textbf{Installation:} \texttt{pip install "pytest>=8"}.

%% ----------------------------------------------------------------------------
\section{Chapter Summary \& Key Takeaways}
\label{sec:ch11-summary}

\begin{itemize}
  \item \textbf{Pattern before protocol.} Choose by initiation and direction:
        server-to-server push $\Rightarrow$ webhooks; one-way server-to-client
        $\Rightarrow$ SSE; genuine two-way interaction $\Rightarrow$
        WebSockets; everything else $\Rightarrow$ request/response, possibly
        with long polling or \texttt{ETag}-assisted short polling.
  \item \textbf{Webhooks are reverse APIs.} At-least-once delivery is the
        only honest semantics; the contract is signed HMAC over
        timestamp-plus-body, a tight tolerance window, bounded jittered
        retries that distinguish permanent from transient failure, idempotent
        receivers, and a dead-letter queue with a replay path.
  \item \textbf{SSE is plain HTTP with a grammar.} Four fields, comment
        keep-alives, blank-line dispatch; \texttt{EventSource} reconnects and
        resumes via \texttt{Last-Event-ID} for free. Respect buffering
        proxies (\texttt{X-Accel-Buffering: no}) and the HTTP/1.1
        connection cap (prefer HTTP/2).
  \item \textbf{WebSocket abandons HTTP after the handshake.} The accept
        value is $\mathrm{Base64}(\mathrm{SHA\text{-}1}(key \,\|\,
        \mathrm{GUID}))$; frames carry FIN/opcode/mask/length; client masking
        protects intermediaries, not secrecy; ping/pong and the close
        handshake are your liveness and shutdown machinery. Do not pay those
        costs for stateless CRUD or cacheable reads.
  \item \textbf{Contracts govern events too.} AsyncAPI documents channels,
        operations, messages, and bindings; compatibility gates and
        idempotent, poison-tolerant consumers are as mandatory as for REST.
  \item \textbf{Test the failure modes, deterministically.} Injected clocks
        and sleepers, scripted transports, \texttt{asyncio.Event}
        synchronization --- never wall-clock assertions --- make tamper,
        replay, retry-storm, and resume behavior provable offline.
\end{itemize}

%% ----------------------------------------------------------------------------
\section{Quick-Reference Card}
\label{sec:ch11-quickref}

\subsection*{Pattern selection matrix (pocket version)}
\begin{table}[h]
  \centering
  \small
  \begin{tabularx}{\textwidth}{l l X}
    \toprule
    \textbf{Need} & \textbf{Pick} & \textbf{Remember} \\
    \midrule
    Server $\to$ server event & Webhook & sign, retry-bounded, DLQ, idempotent receiver \\
    Server $\to$ browser, one-way & SSE & resume via \texttt{Last-Event-ID}; kill proxy buffering \\
    Two-way low-latency & WebSocket & own reconnection, heartbeat, backpressure \\
    Occasional freshness, cacheable & Short poll + \texttt{ETag} & nearly free at scale \\
    Simple push through hostile proxies & Long polling & hold deadline $<$ proxy idle timeout \\
    \bottomrule
  \end{tabularx}
\end{table}

\subsection*{SSE field reference}
\begin{table}[h]
  \centering
  \small
  \begin{tabularx}{\textwidth}{l X}
    \toprule
    \textbf{Line form} & \textbf{Meaning} \\
    \midrule
    \texttt{data: <text>} & payload; repeated lines joined with newlines \\
    \texttt{event: <name>} & dispatch type (default \texttt{message}) \\
    \texttt{id: <text>} & sets the client's last-event-id cursor \\
    \texttt{retry: <ms>} & reconnection delay (digits only) \\
    \texttt{: <anything>} & comment; use as keep-alive every 10--15\,s \\
    (blank line) & dispatch the buffered event \\
    \bottomrule
  \end{tabularx}
\end{table}

\subsection*{WebSocket close codes}
\begin{table}[h]
  \centering
  \small
  \begin{tabularx}{\textwidth}{l l X}
    \toprule
    \textbf{Code} & \textbf{Name} & \textbf{Use} \\
    \midrule
    1000 & normal & work complete; default clean close \\
    1001 & going away & server restart, page navigation \\
    1002 & protocol error & peer violated framing rules \\
    1003 & unsupported data & received a type you cannot accept \\
    1006 & (abnormal) & reserved; never sent on the wire \\
    1007 & invalid payload & e.g.\ non-UTF-8 in a text frame \\
    1008 & policy violation & authz failure, generic policy refusal \\
    1009 & too big & message exceeded your limit \\
    1011 & internal error & server-side unexpected condition \\
    \bottomrule
  \end{tabularx}
\end{table}

\subsection*{Webhook header anatomy}
\begin{minted}{text}
X-Northwind-Signature: t=<unix-seconds>,v1=<hex HMAC-SHA256(secret, "t.body")>
verify order: parse -> freshness (|now - t| <= 300 s) -> HMAC (constant time)
              -> parse JSON -> idempotency check -> process -> 200
retry policy: 2xx = done; 408/409/425/429/5xx/network = retry (bounded,
              jittered exponential); other 4xx = dead-letter immediately
\end{minted}

\section{Standardized Evidence Addendum}
\subsection{Benchmark Snapshot Template}
\begin{table}[h]
  \centering
  \small
  \begin{tabularx}{\textwidth}{l c c c c}
    \toprule
    \textbf{Config} & \textbf{p50 latency} & \textbf{p99 latency} & \textbf{Error rate} & \textbf{Throughput} \\
    \midrule
    Baseline & -- & -- & -- & 1.00x \\
    Candidate A & -- & -- & -- & -- \\
    Candidate B & -- & -- & -- & -- \\
    \bottomrule
  \end{tabularx}
  \caption{Minimum benchmark table to keep per chapter.}
\end{table}

\subsection{Request Trace Template}
\begin{itemize}
  \item \textbf{Request:} one representative API call for this chapter's topic.
  \item \textbf{Before:} behavior (headers, payload, latency, failure mode) with the naive approach.
  \item \textbf{After:} behavior with the technique in this chapter.
  \item \textbf{Result:} what changed in correctness, resilience, latency, and operability.
\end{itemize}

\subsection{Cost and Latency Budget Snapshot}
\begin{table}[h]
  \centering
  \small
  \begin{tabularx}{\textwidth}{l c c}
    \toprule
    \textbf{Stage} & \textbf{Latency budget} & \textbf{Cost driver} \\
    \midrule
    TLS / connection setup & -- & handshake round trips, keep-alive hit rate \\
    Request validation & -- & schema complexity, CPU per request \\
    Business logic and I/O & -- & downstream fan-out, query cost \\
    Serialization and egress & -- & payload size, compression ratio \\
    \bottomrule
  \end{tabularx}
  \caption{Operational budget template for this chapter's technique.}
\end{table}

\section{Configuration Preset Template}
\begin{minted}{text}
# api_policy.yaml
policy_version: "v1.0.0"
component: "<chapter_topic>"
timeout_ms: 0
retry_budget: 0
fallback_strategy: "fail_fast"
rollback_trigger: "error_budget_burn_gt_threshold"
\end{minted}

\section{CI Quality Gate Specification}
\begin{itemize}
  \item contract tests green against the pinned OpenAPI/AsyncAPI/Protobuf spec,
  \item no breaking schema changes without a version bump,
  \item error responses conform to RFC 9457 problem-details shape,
  \item benchmark deltas within approved regression bounds (latency and throughput),
  \item policy version and rollback plan present in release metadata.
\end{itemize}

\section{Incident Runbook Template}
\begin{enumerate}
  \item Detect regression from SLO burn-rate alerts or consumer reports.
  \item Scope impacted endpoints, versions, and consumer segments.
  \item Contain impact (rollback, feature flag, rate-limit tightening, or traffic shift).
  \item Diagnose with traces, structured logs, and contract diffs.
  \item Recover with validated fix and verified rollout procedure.
  \item Prevent recurrence via new regression tests and CI gates.
\end{enumerate}

\section{Cross-Chapter Decision Card}
\begin{table}[h]
  \centering
  \small
  \begin{tabularx}{\textwidth}{l X X}
    \toprule
    \textbf{Dimension} & \textbf{Choose this approach when} & \textbf{Prefer an alternative when} \\
    \midrule
    Use case & -- & -- \\
    Latency budget & -- & -- \\
    Operational cost & -- & -- \\
    Team maturity & -- & -- \\
    \bottomrule
  \end{tabularx}
  \caption{Decision card to complete per chapter.}
\end{table}

\section{ROI Model and Build-vs-Buy}
\begin{itemize}
  \item \textbf{Build when:} the capability is differentiating, compliance forbids vendors, or scale makes vendor pricing irrational.
  \item \textbf{Buy when:} the capability is commodity (auth, gateways, observability), time-to-market dominates, or staffing is thin.
  \item \textbf{Hidden costs to model:} on-call load, upgrade cadence, expertise ramp, data egress fees.
\end{itemize}

\section{Disaster Recovery Notes (RPO/RTO)}
\begin{itemize}
  \item Define recovery point objective (RPO): maximum acceptable data loss window.
  \item Define recovery time objective (RTO): maximum acceptable restoration time.
  \item Test restores regularly; an untested backup is not a backup.
\end{itemize}

